\documentclass[journal]{IEEEtran}
\usepackage{amsmath,amssymb,mathtools}
\usepackage{amsthm}
\usepackage{booktabs,array}
\usepackage{enumitem}
\usepackage{microtype}
\usepackage[hidelinks]{hyperref}
\newtheoremstyle{plainn}{}{}{\itshape}{}{\bfseries}{.}{.5em}%
  {\thmname{#1}\thmnumber{ #2}\thmnote{ (#3)}}
\theoremstyle{plainn}

\newtheorem{thm}{Theorem}
\newtheorem{prop}{Proposition}
\newtheorem{cor}{Corollary}
\newtheorem{claim}{Claim}
\newtheorem{defn}{Definition}
\newtheorem{exmp}{Example}
\newtheorem{algo}{Algorithm}
\theoremstyle{remark}
\newtheorem{remark}{Remark}
\renewcommand{\qedsymbol}{QED}
\newcommand{\row}{\operatorname{row}}
\newcommand{\rank}{\operatorname{rank}}
\newcommand{\Ann}{\operatorname{Ann}}
\newcommand{\Gr}{\operatorname{Gr}}
\newcommand{\Prob}{\operatorname{Prob}}
\newcommand{\poly}{\operatorname{poly}}

\begin{document}

\title{Functional-Observer Feasibility in the Hyperplane Layer}
\author{}
\maketitle

\begin{abstract}
The companion papers decide whether a functional observer of a prescribed order exists and compute the
minimum order; neither describes the admissible designs. This paper does, for the smallest layer in
which the kernel of the compressed pencil is one-dimensional. The admissible kernel directions are
exactly the real eigenvectors of the orthogonal compression of the state matrix to the unobservable
subspace, further restricted by the prescribed rows. The design family is therefore a finite union of
relative Grassmann bundles over locally closed algebraic strata of projective eigenspaces, with
explicit fibre dimensions and nonemptiness conditions. Feasibility in the layer is decided exactly by
algebraic computation on these bundles, and a full-rank leading-coefficient lemma lets a single
computation per chart settle every point of it rather than a generic one. For fixed output dimension
and geometrically simple real spectrum the decision is polynomial in the input bit size. A certified
repeated-eigenspace example shows that feasibility need not be constant on a projective eigenspace, so
one eigenvector per eigenvalue is not a valid reduction; no hardness statement is claimed. A duality
identifies the projective direction with the annihilator of the common projection of the compressed
formulation.
\end{abstract}

\begin{IEEEkeywords}
Functional observers, feasibility, compressed pencil, Grassmannian, exact algorithms,
parameterized complexity.
\end{IEEEkeywords}

\section{Introduction}\label{sec:intro}

For a linear system $(A,C)$ with a prescribed block $L_0$ of functional rows, a design is a matrix
$\mathcal L=[L_0;R]$ with $q$ independent rows, and it is feasible when
\begin{equation}\label{eq:5a}
\rank\begin{bmatrix}\mathcal LA\\ \mathcal L\\ CA\\ C\end{bmatrix}
=\rank\begin{bmatrix}\mathcal L\\ CA\\ C\end{bmatrix},
\end{equation}
\begin{equation}\label{eq:5b}
\rank\begin{bmatrix}\lambda\mathcal L-\mathcal LA\\ CA\\ C\end{bmatrix}
=\rank\begin{bmatrix}\mathcal L\\ CA\\ C\end{bmatrix}\ \ \forall\lambda\in\mathbb C .
\end{equation}
The questions one usually asks about \eqref{eq:5a}--\eqref{eq:5b} are settled by the companion
papers; what this paper asks is a different one.

The first \cite{companion1} decides feasibility for a prescribed order, after compressing the free
rows to a subspace of a compressed row-pair space; the second \cite{companion2} computes the minimum
order and constructs a minimum-order observer from a subspace $W$ of that space, with the structural
bound $\omega+1\le\dim W\le 2\omega$. Neither asks what the designs themselves look like. This paper
describes them for the smallest layer in which a direction exists, that is, when the kernel of the
compressed pencil is one-dimensional. We call the set of all admissible designs of that layer its
\emph{algebraic parameter space}, and we describe it as a finite union of fibre bundles; each bundle
sits over a \emph{locally closed algebraic stratum} of a projective eigenspace, and its fibres are
Grassmannians. No single closed subvariety describes the whole parameter space, and the non-eigen
stratum in particular is open rather than closed, so the stratum-by-stratum statement is the natural
one.

The results are the following.

\begin{enumerate}[label=(\Alph*),leftmargin=*]
\item \emph{The kernel direction is an eigenvector.} The kernel line of $H=\mathcal LN$ is not free: it
is generated by a real eigenvector $v=Na$ of the orthogonal compression $M=\pi_{K_0}A|_{K_0}$, where
$a\in\mathbb R^k$ collects the coordinates of the direction in the basis $N$ of $K_0$. With prescribed
rows the admissible directions are further restricted to
$E_\mu^0=E_\mu\cap\ker(L_0N)\cap\ker(L_0AN)$ (Theorem~\ref{sec3:thm:direction},
Claim~\ref{sec3:lem:prescribed}). This turns "choose $q$ rows" into "choose a direction, then a
subspace over it".
\item \emph{The algebraic parameter space is a finite union of relative Grassmann bundles.} Over each
locally closed stratum of each projective eigenspace --- the eigen directions and the non-eigen
directions --- the admissible designs form a relative Grassmann bundle whose fibre dimension and
nonemptiness condition are explicit (Theorem~\ref{sec3:thm:param}). The fibre dimension is
$(s-\delta)(r(v)-s)$, where $s$ is the number of rows beyond the compulsory $k-1$, $r(v)\in\{\rho,\rho-1\}$
is the rank of the quotient bundle, and $\delta$ measures the prescribed rows modulo $K_0$.
\item \emph{One computation per chart, valid at every point of it.} On a chart of the parameter space,
the maximal minors have no common root over the function field if and only if \emph{every} point of
that chart is feasible. The step from a generic point to the whole chart is a full-rank
leading-coefficient lemma (Claim~\ref{sec4:lem:fullrank}); without it a generic test certifies only a
generic point. This is what makes the chartwise computation a decision procedure rather than a
sampling heuristic, and it is the property that the counterexample below would otherwise violate.
\item \emph{Fixed-$p$ tractability, and where it stops.} For fixed output dimension $p$ and a
geometrically simple real spectrum of $M$, the decision is polynomial in the input bit size
(Theorem~\ref{sec5:thm:poly}), with the exponent depending on $p^2$. The hypothesis governs the
dimension of the direction parameter, and it cannot simply be dropped; no hardness statement for the
repeated case is claimed.
\item \emph{A certified counterexample.} Inside a two-dimensional eigenspace one direction is feasible
and five others are not, with complete rational data and exact certificates
(Appendix~\ref{app:example}). Hence "one eigenvector per eigenvalue" is not a valid reduction.
\item \emph{A duality with the compressed formulation.} In this layer the projective direction is the
annihilator of the common projection used in \cite{companion2}: with $U=\pi_1(W)=\pi_2(W)$ one has
$U=a^\circ$ and $[a]=\mathbb P(\Ann U)$ (Proposition~\ref{sec3:prop:duality}).
\end{enumerate}

Three statements of different strength are made, and they are kept apart throughout: (i) an exact
decision procedure for the whole layer, in exact arithmetic and with no complexity claim
(Section~\ref{sec4}); (ii) a polynomial bound on the \emph{cost} of that same procedure on the
geometrically simple subclass, a promise-class statement about running time rather than about which
inputs are decided (Section~\ref{sec5}); (iii) a counterexample refuting one specific reduction
(Appendix~\ref{app:example}). Nothing is claimed about the complexity of the repeated case.


The paper is organized as follows: Section~\ref{sec:prelim} fixes what is borrowed and the notation,
Section~\ref{sec3} describes the layer, Section~\ref{sec4} decides feasibility on it,
Section~\ref{sec5} bounds the cost, Section~\ref{sec6:sub:conclusion} states what is not claimed,
Appendix~\ref{app:example} holds the counterexample and Appendix~\ref{appBC} the proofs.

\section{Standing results and notation}\label{sec:prelim}

This section fixes the ground the rest of the paper stands on: what is taken from the companion
papers without proof, and what notation is specific to the hyperplane layer.

\subsection{Results taken from Paper I}\label{sec:prelim:criterion}

Three things are taken from \cite{companion1} and not reproved here: the decision problem with its
fixed-order answer, the compressed criterion, and the certificate built for the spectral condition.

\emph{The problem.} $\mathrm{DEC}$ takes $(A,C,L_0)$ and $q$ and asks whether some $R$ satisfies
\eqref{eq:5a}--\eqref{eq:5b}; no witness is required. Theorem~1 of \cite{companion1} answers it in
deterministic polynomial time for every fixed $q$, in $\mathcal L^{O(q^2)}$ bit operations for input
bit length $\mathcal L$; Theorem~2 there parameterizes instead by $\dim S$ and pays an exponential in
it. The parameterization of Section~\ref{sec5} fixes the output dimension $p$, so its bit bound is
proved here.

\emph{The criterion (Claim~2 of \cite{companion1}).} In the compressed data the conditions read as
follows.

The matrix $\mathcal L$ is feasible, that is, it satisfies \eqref{eq:5a} and \eqref{eq:5b}, if and only
if
\begin{equation}
  \label{eq:5ap}
  \row H_A\subseteq\row H,
\end{equation}
\begin{equation}
  \label{eq:5bp}
  \rank(\lambda H-H_A)=\omega\qquad\text{for all }\lambda\in\mathbb C .
\end{equation}
This equivalence, the identity $\rank[\mathcal L;CA;C]=\rank[CA;C]+\omega$, and the reformulation of
\eqref{eq:5bp} as the absence of a common root of the maximal minors of the pencil are established in
\cite{companion1}; they are used here and not re-proved. Condition \eqref{eq:5ap} only confines the
problem to a parameter space. What this paper adds to the criterion is a statement about \emph{which
points} a generic test covers, which is Claim~\ref{sec4:lem:fullrank}.


\emph{The certificate (Claims~8 and 11 of \cite{companion1}).} That paper replaces the maximal minors
by a pivot factorization, a rank-consistency test (Claim~8) and its Gram-matrix compression
(Claim~11), of a size governed by $\omega\le\ell$. This paper forms the minors instead --- in the layer
there are few of them --- but keeps the reading of the criterion on which that certificate rests.

The purpose of this step is to fix which statements used here are standard and which are proved here; it
is used in Section~\ref{sec4}.

\begin{claim}[Standard tools: bad locus, generic point, finite cover]
\label{sec4:thm:closed}
\label{sec4:thm:generic}
\label{sec4:prop:cover}
Let the layer be covered by finitely many irreducible parameter pieces with $\rank H\equiv\omega$ and
function field $K=\overline{\mathbb Q(\mathcal S)}$. Then:
\begin{enumerate}[label=(\alph*),leftmargin=*]
\item \label{it:closed} for every chart $\mathcal X$ the bad locus $\mathcal B\subseteq\mathcal X$ is
Zariski closed;
\item \label{it:generic} a piece $\mathcal S$ contains a good point if and only if
$\deg_\lambda\gcd_I m_I=0$ in $K[\lambda]$;
\item \label{it:cover} consequently $\mathcal X$ contains a good point if and only if some piece of the
cover has a gcd of $\lambda$-degree zero.
\end{enumerate}
\end{claim}


These statements are elementary facts about a family of polynomials over an irreducible base; none of them
is taken from \cite{companion1}, which contains no closedness result, and the parts that are not immediate
are proved in this paper. Statement (a) is
the closedness of the projection to $\mathcal X$ of the incidence set
$\{(X,(\lambda:\mu)):\tilde m_I(X,\lambda,\mu)=0\ \forall I\}\subseteq\mathcal X\times\mathbb P^1$, and it
is proved for the charts of the layer in Appendix~\ref{appBC}, paragraph \emph{Why the bad locus is
closed}. Statement (c) is immediate from (b) and the finiteness of the cover.
Statement (b) splits: the implication from $\deg_\lambda\gcd=0$ to the existence of a good
point is immediate, since a gcd of $\lambda$-degree zero is a unit; the converse --- a good point forces
$\deg_\lambda\gcd=0$ --- is Claim~\ref{sec4:lem:fullrank} below, and it is proved there from
$\rank H\equiv\omega$. For the general statements on gcds over a function field see \cite{lan02,vzgg03}.

We follow Sections~\ref{sec:prelim} and~\ref{sec3}. Let $\mathcal L=\binom{L_0}{R}$ have $q$ rows and
$V=\row\mathcal L$, let $V_0=\row L_0$, $N$ be a basis of $K_0$, and
\[
\begin{gathered}
H=LN,\qquad H_A=LAN,\\
P(\lambda)=\lambda H-H_A,\qquad \omega=\rank H ,
\end{gathered}
\]
with $\omega=k-1$ in the hyperplane layer, $\rho=\rank[C;CA]$, $s=q-k+1$, and
$\widetilde r=r(v)-\delta$ the rank of the quotient bundle. A parameter point $X$ of a chart gives
$L(X)$, $H(X)=L(X)N$, $H_A(X)=L(X)AN$ and $P(X,\lambda)=\lambda H(X)-H_A(X)$, a $q\times k$ pencil of
rank $\omega$ at every $X$. The index $I$ ranges over pairs of $\omega$ rows and $\omega$ columns, and
$m_I(X,\lambda)$ is the corresponding maximal minor.


\subsection{Results taken from Paper II}\label{sec:prelim:assumptions}

The second companion paper \cite{companion2} works with $W=\row[H\ H_A]\subseteq\mathbb R^{2k}$, the
image of a row space under $\mathsf M=[N\ \ AN]$, and with the projections $\pi_1(h,g)=h$,
$\pi_2(h,g)=g$ and $\pi_\lambda(h,g)=\lambda h-g$. Its Claim~3 reads the two conditions on $W$ alone:
\eqref{eq:5a} holds if and only if $\pi_2(W)\subseteq\pi_1(W)$, and \eqref{eq:5b} if and only if
$\dim\pi_\lambda(W)=\dim\pi_1(W)$ for all $\lambda$, every such $W$ being realizable by a free block
of $t=\dim(W/W_0)$ rows. Its Claim~4 is the structural bound: for a feasible $W$ with
$\omega=\dim\pi_1(W)>0$ one has $\omega+1\le\dim W\le 2\omega$, and $\pi_1(W)=\pi_2(W)=\row H_A$.

For orientation: Claim~5 there gives monotonicity and saturation of the feasible orders, Proposition~1
that $q_{\min}=r_0+t_{\min}$ with $t_{\min}\le2\ell$, Proposition~2 a lower bound
$r_0+\max(0,\operatorname{rank}[U_0;V_0]-s_0+1)$, and Proposition~3 when that bound is attained. None
of these is used below; what is used is the subspace reading and the structural bound, in the duality
at the end of Section~\ref{sec3}.

\subsection{Notation specific to this paper}\label{sec:prelim:notation}

Three blocks follow: the system and its block structure, the direction and the two functionals, and
the parameter counts with the chart conventions.

\emph{The system and its block structure.}

Let $A\in\mathbb R^{n\times n}$, $C\in\mathbb R^{p\times n}$ and $L_0\in\mathbb R^{r_0\times n}$. The
design matrix is $\mathcal L=[L_0;R]\in\mathbb R^{q\times n}$, assumed to have full row rank, and its
row space is $V=\row\mathcal L\subseteq\mathbb R^n$. Put
\begin{equation}
  \label{eq:K0}
  K_0=\ker\left[\begin{smallmatrix}C\\ CA\end{smallmatrix}\right],\qquad
  k=\dim K_0,\qquad
  \rho=n-k,
\end{equation}
where $\rho=\rank[C;CA]$. Fix a basis $N\in\mathbb R^{n\times k}$ of $K_0$; Appendix~\ref{appB}
describes a normalized
choice. Here $\rho$ is the rank of the compressed data, whereas the layer index of this paper is $c$;
the two are different quantities. For a row space $V$ put
\begin{equation}
  \label{eq:compressed-data}
  H=LN,\qquad H_A=LAN,\qquad \omega=\rank H,\qquad c=k-\omega .
\end{equation}
The \emph{layer} is $c$. This paper treats only $c=1$, that is, $\omega=k-1$.


The conventions that hold throughout are the following; the third is forced by the block structure
rather than assumed.

Throughout the paper:
\begin{enumerate}[label=(\roman*),leftmargin=*,itemsep=2pt]
\item the standard Euclidean inner product is used, so restrictions to subspaces are automatically
nondegenerate;
\item input matrices are rational;
\item the design matrix $\mathcal L$ has full row rank, hence $\rank L_0=r_0$; otherwise no admissible
$\mathcal L$ exists at all, see Claim~\ref{prop:rankprecheck};
\item the parameterization requires a basis of $K_0$; if $N'=NT$ with $T\in\operatorname{GL}_k(\mathbb R)$
is another basis, then $H'=HT$ and $H_A'=H_AT$, so $\lambda H'-H_A'=(\lambda H-H_A)T$. The change is
\emph{right} multiplication by an invertible matrix, which preserves ranks and hence feasibility.
\end{enumerate}


The purpose of this step is to record that $\rank L_0=r_0$ is forced rather than assumed; it is used
in Algorithm~\ref{sec4:algo:decide}.

\begin{claim}
  \label{prop:rankprecheck}
  If $\rank L_0<r_0$, then for every $q$ there is no matrix $\mathcal L=[L_0;R]$ of full row rank, so
  the instance is infeasible.
\end{claim}

\begin{proof}
  Such a matrix would satisfy
  $q=\rank\mathcal L\le\rank L_0+(q-r_0)<r_0+(q-r_0)=q$, a contradiction.
\end{proof}


\emph{The direction and the two functionals.} The direction is where the layer differs from the
compressed formulation, so it receives notation of its own.

Let $V$ lie in the hyperplane layer and assume the layer transverse, that is
$\dim(V\cap K_0)=k-c=k-1$; this is hypothesis \ref{sec3:S2} of Section~\ref{sec3}, not a consequence of
\eqref{eq:compressed-data} (Remark~\ref{sec3:rem:identity}). By \eqref{eq:compressed-data} and $c=1$ one
has $\omega=k-1$, hence $\ker H$ is a line $\mathbb R v$ with $v\in K_0\setminus\{0\}$, and
$V\subseteq U_v^\perp$, where
\begin{equation}
  \label{eq:Uv}
  U_v=\operatorname{span}\{v,Av\}.
\end{equation}
Denote the two linear functionals on $K_0$ by
\begin{equation}
  \label{eq:functionals}
  \varphi_v(u)=\langle u,v\rangle,\qquad \varphi_{Av}(u)=\langle u,Av\rangle
  \qquad(u\in K_0).
\end{equation}
The inclusion $V\subseteq U_v^\perp$ is equivalent to saying that every $x\in V$ annihilates both $v$
and $Av$, that is, that $xN$ lies in the common kernel of the two functionals.

\begin{remark}[$\omega$ is not $\dim(V\cap K_0)$]\label{sec3:rem:identity}
By \eqref{eq:compressed-data} and $\ker(L|_{K_0})=K_0\cap V^\perp$ one has
$\omega=\rank H=k-\dim(K_0\cap V^\perp)=\dim\pi_{K_0}(V)$, with $\pi_{K_0}$ the orthogonal projection
onto $K_0$; so the layer index is
$c=k-\omega=\dim(K_0\cap V^\perp)$ and counts the codimension of $\pi_{K_0}(V)$ in $K_0$, not the
intersection $V\cap K_0$. The naive identity $\omega=\dim(V\cap K_0)$ is false: an instance with $n=8$,
$k=4$ and a design of $q=5$ rows has $\omega=\rank H=4=k-\dim(K_0\cap V^\perp)$, while
$\dim(V\cap K_0)=2$. The two numbers agree exactly when $\dim(V\cap K_0)=k-c$, which is the
transversality hypothesis \ref{sec3:S2}.
 \emph{Scope.} The parameterization below covers the designs with $V\supseteq D_v$ (equivalently $\dim(V\cap K_0)=k-c$). Whether every feasible design of the layer admits a feasible design with this property, at the same order, we leave open; all decision results in this paper are therefore stated for that subfamily. The statements are given projectively; the formalization uses the equivalent linear (cone) version, the two differing only in notation.
\end{remark}

\begin{equation}
  \label{eq:compression}
  \begin{aligned}
    M&=\pi_{K_0}A|_{K_0},\qquad A_{21}=\pi_{K_0^\perp}A|_{K_0},\\
    E_\mu&=\ker(M-\mu I),\\
    E_\mu^0&=E_\mu\cap\ker(L_0N)\cap\ker(L_0AN),
  \end{aligned}
\end{equation}
\begin{equation}
  \label{eq:fibre-data}
  \begin{aligned}
    D_v&=K_0\cap U_v^\perp,\qquad \mathcal Q_v=U_v^\perp/D_v,\\
    V_0&=\row L_0,\qquad s=q-k+1,\\
    \delta&=\dim V_0-\dim(V_0\cap K_0),
  \end{aligned}
\end{equation}
together with the rank $r(v)=\dim\mathcal Q_v$ of the fibre. Its two values are computed in
Section~\ref{sec3}: $r(v)=\rho$ for eigen directions and $r(v)=\rho-1$ for non-eigen
directions.

The standard Euclidean inner product on $\mathbb R^n$ is used throughout, so restrictions to subspaces
are nondegenerate; inputs are rational and $\mathcal L$ has independent rows. Table~\ref{tab:notation}
collects the symbols; it is repeated nowhere else and is used as the reference for the whole paper.
Two points of the table deserve emphasis. First, $v$ is an ambient vector in $K_0\subseteq\mathbb R^n$
while $a\in\mathbb R^k$ is its coordinate vector, $v=Na$; the kernel direction is a line in the first
space and a projective point in the second, and it is the second that appears in the duality of
contribution~(F). Second, $\rho=\rank[C;CA]$ is the rank of the compressed data and is not the layer
index, which is $c=k-\omega$.

\begin{table*}[t]
\centering
\caption{Symbols used throughout, in two blocks read row by row. The third cell of each block gives
the space and its dimension; ``---'' marks a scalar or an index.}
\label{tab:notation}
\small
\begin{tabular}{@{}llllll@{}}
\toprule
Symbol & Object & Space, dimension & Symbol & Object & Space, dimension \\
\midrule
$K_0$ & $\ker[C;CA]$, unobservable subspace & $\mathbb R^n$, $k$ &
$E_\mu$ & $\ker(M-\mu I)$, $\mu$ real & $K_0$, $d_\mu$ \\
$k$ & $\dim K_0$ & --- &
$E_\mu^0$ & $E_\mu\cap\ker(L_0N)\cap\ker(L_0AN)$ & $K_0$, $d_\mu^0$ \\
$\rho$ & $\rank[C;CA]=n-k$ & --- &
$v$, $a$ & direction and coordinates, $v=Na$ & $K_0$ resp.\ $\mathbb R^k$ \\
$N$ & basis matrix of $K_0$ & $\mathbb R^{n\times k}$ &
$U_v$ & $\operatorname{span}\{v,Av\}$ & $\mathbb R^n$, $1$ or $2$ \\
$H$, $H_A$ & $\mathcal LN$, $\mathcal LAN$ & $\mathbb R^{q\times k}$ &
$D_v$ & $K_0\cap U_v^\perp$ & $K_0$, $k-1$ \\
$\omega$ & $\rank H=k-\dim(K_0\cap V^{\perp})$ & --- &
$\mathcal Q_v$ & $U_v^\perp/D_v$ & quotient, $r(v)$ \\
$c$ & layer index $k-\omega$ & --- &
$V_0$ & $\row L_0$, prescribed rows & $\mathbb R^n$, $r_0$ \\
$M$ & $\pi_{K_0}A|_{K_0}$ & $\operatorname{End}K_0$, $k\times k$ &
$\delta$ & $\dim V_0-\dim(V_0\cap K_0)$ & --- \\
$A_{21}$ & $\pi_{K_0^\perp}A|_{K_0}$ & $K_0\to K_0^\perp$ &
$s$ & $q-k+1$, rows beyond $k-1$ & --- \\
$r(v)$ & $\rho$ (eigen), $\rho-1$ (non-eigen) & --- &
$\widetilde r$ & $r(v)-\delta$, rank left for free rows & --- \\
\bottomrule
\end{tabular}
\end{table*}

The eigen/non-eigen stratification (Proposition~\ref{sec3:prop:strata}) is formalised in \texttt{Strata.lean} (41 declarations, abstract core); the projectivisation and the Zariski-open statement are paper-level, the Lean version uses the equivalent linear (cone) form.
The specialisation half of Lemma~\ref{sec4:lem:fullrank} is formalised as \texttt{specialized\_common\_factor} (with \texttt{exists\_common\_factor\_of\_gcd} for a family of points and \texttt{exists\_eval\_ne\_zero} for the nonvanishing half); the chart layer and the Zariski-open statement \texttt{sec4:thm:generic} remain paper-level.


\emph{Parameter counts and chart conventions.} The last block fixes what a parameter point is and over
which field it is computed.

The parameter space is taken to be a finite union of \emph{affine charts}: affine charts of projective
spaces for the direction, affine charts of Grassmannians for the fibre, and -- when
$\dim E_\mu^0>1$ -- affine charts of the direction base as well, combined with local trivializations of
the quotient bundle. For a real algebraic eigenvalue $\mu$ these charts are constructed over the number
field $K=\mathbb Q(\mu)$ equipped with the real embedding that selects $\mu$; their scalar extensions to
$\mathbb R$ are affine spaces whose real points are Zariski dense. All chart polynomials and
function-field gcd computations are performed over $K$. This convention is used in
Section~\ref{sec4} to pass from a generic algebraic point to a realizable real design.


\section{The hyperplane layer}\label{sec3}

This section describes the layer: which directions the compressed pencil admits, which of those
survive the prescribed rows, and what the admissible row spaces over a fixed direction look like. The
outcome is an explicit finite union of relative Grassmann bundles.

The admissible designs form a finite union of \emph{relative Grassmann bundles}, with bases the real
projective eigenspaces of the compression restricted by the prescribed rows and explicit Grassmannian
fibres; feasibility in the layer is thereby reduced to an algebraic decision problem on finitely many
parameter spaces (Section~\ref{sec4}). The inner product is the standard Euclidean one on
$\mathbb{R}^n$ \cite{hj13}.

The parameterization has two independent steps. First the \emph{direction}: the kernel direction must
be an eigenvector of the compression (Claim~\ref{sec3:lem:direction} and
Theorem~\ref{sec3:thm:direction}), and the prescribed rows restrict it to a linear subspace of that
eigenspace (Claim~\ref{sec3:lem:prescribed}). Then the \emph{fibre} over a fixed direction: the
admissible row spaces form a Grassmannian whose rank is read off from the eigen/non-eigen
stratification (Proposition~\ref{sec3:prop:strata} and Claim~\ref{sec3:lem:constantrank}).


\subsection{The layer, its direction, and the direction theorem}\label{sec3:sub:struct}\label{sec3:sub:direction}

Section~\ref{sec:prelim} fixes $K_0$, $k$, $\rho$, $N$, $H$, $H_A$, $\omega$, $c$, $M$, $A_{21}$,
$E_\mu$, $U_v$, $D_v$, $\mathcal Q_v$, $V_0$, $\delta$, $s$. Five symbols are added:
$\operatorname{Spec}_{\mathbb R}(M)$ for the set of real eigenvalues of $M$,
$E_\mu^0=E_\mu\cap\ker(L_0N)\cap\ker(L_0AN)$,
$h_\mu^0=\dim(E_\mu^0\cap\ker A_{21})$, $d_\mu^0=\dim E_\mu^0$, and $\Gr(m,Z)$ for the Grassmannian of
$m$-dimensional subspaces of a vector space $Z$.

\noindent\emph{Standing assumption.} $\rank L_0=r_0$, which is forced by feasibility
(Claim~\ref{sec3:proc}).

A \emph{row space of the layer} is a $q$-dimensional subspace $V$ satisfying
\begin{enumerate}[label=(S\arabic*),ref={(S\arabic*)},leftmargin=2.6em,itemsep=1pt,topsep=2pt]
\item\label{sec3:S1} $V\subseteq U_v^\perp$, where $v\in K_0\setminus\{0\}$ is determined by $V$;
\item\label{sec3:S2} $D_v\subseteq V$ (given \ref{sec3:S1}, equivalently $\dim(V\cap K_0)=k-1=k-c$);
\item\label{sec3:S3} $V_0\subseteq V$.
\end{enumerate}
Condition~\ref{sec3:S3} belongs to the problem definition, and \ref{sec3:S1} with \ref{sec3:S2} gives
$V\cap K_0=D_v$. Condition~\ref{sec3:S2} is an explicit hypothesis on the layer, not a consequence of
\eqref{eq:5a}: with $M=\operatorname{diag}(1,1,2,3)$, $K_0=\operatorname{span}\{e_5,\dots,e_8\}$,
$D=\operatorname{span}\{e_5+e_6,\ e_7+e_8\}\subseteq K_0$ not $M$-invariant, $U=D+AD$ and $V=U^\perp$ one
has $c=2$ and \eqref{eq:5a}, but $\dim(V\cap K_0)=1\ne k-c$; what \eqref{eq:5a} yields in general is only
the inequality $\dim(V\cap K_0)\le\omega$ of Claim~\ref{sec3:lem:direction}(iv).


The purpose of this step is to name the direction and the space $D_v$ that every later count is built
from; it is used in Theorem~\ref{sec3:thm:direction}.

\begin{claim}[Structure of the layer]
\label{sec3:lem:direction}
\label{sec3:lem:Dv}
\label{sec3:rem:omega}
Let $\mathcal{L}$ be a design with $c=1$ satisfying \eqref{eq:5a}, and write
$\varphi_v(u)=\langle u,v\rangle$, $\varphi_{Av}(u)=\langle u,Av\rangle$ for the two functionals on
$K_0$. Then:
\begin{enumerate}[label=(\roman*),leftmargin=*,itemsep=1pt,topsep=2pt]
\item $\ker(\mathcal{L}|_{K_0})=K_0\cap V^\perp=\mathbb{R} v$ for a unique
$v\in K_0\setminus\{0\}$, determined up to a nonzero scalar;
\item $\mathcal{L}v=0$ and $\mathcal{L}Av=0$, that is, $V\subseteq U_v^\perp$; if moreover
$V_0\subseteq V$ then $L_0v=L_0Av=0$;
\item $\dim D_v=k-\rank\{\varphi_v,\varphi_{Av}\}\ge k-2$, and for a design of the layer
$\dim D_v=k-1$;
\item $\dim(V\cap K_0)\le\omega$, with equality if and only if $D_v\subseteq V$; so \ref{sec3:S2} is a
normalization in the definition of a layer design, not a consequence of compression.
\end{enumerate}
\end{claim}
\begin{proof}
Since $\ker H$ is one-dimensional, so is $\{\mathcal{L}u=0,\ u\in K_0\}=N\ker H$; let $v$ generate it,
which gives $\mathcal{L}v=0$, half of (ii). By \eqref{eq:5a} there is $X$ with $H_A=XH$, hence
$H_Ac=XHc=0$ for every $c\in\ker H$; taking $c$ with $Nc=v$ gives $\mathcal{L}Av=0$, and if
$V_0\subseteq V\subseteq U_v^\perp$ every row of $L_0$ lies in $U_v^\perp$, i.e.\ $L_0v=L_0Av=0$.
Note that \eqref{eq:5a} enters only through $\ker H\subseteq\ker H_A$, that is, through
$\row H_A\subseteq\row H$.

For (iii), $K_0\cap U_v^\perp$ is the common kernel of the two functionals, whose dimension is $k$
minus their rank. For (iv), $\omega=\dim\pi_{K_0}(V)$ and $V\cap K_0\subseteq\pi_{K_0}(V)$, so
$\dim(V\cap K_0)\le\omega=k-1$; equality holds if and only if $D_v\subseteq V$, which is
\ref{sec3:S2}. By (iii) a layer design therefore has $\rank\{\varphi_v,\varphi_{Av}\}=1$, the fact used in Claim~\ref{sec3:lem:constantrank}.
\end{proof}


The direction cannot be an arbitrary line of $K_0$, and the next statement says which lines occur as
kernel directions of row spaces satisfying \ref{sec3:S1}--\ref{sec3:S2}; the prescribed rows
\ref{sec3:S3} are a further condition and are treated separately in Claim~\ref{sec3:lem:prescribed}.

\begin{thm}[Direction theorem]
\label{sec3:thm:direction}
\label{sec3:rem:direction-set}
\label{sec3:rem:equiv}
Let $v\in K_0\setminus\{0\}$. Then there is a $q$-dimensional row space $V$ satisfying
\ref{sec3:S1}--\ref{sec3:S2}, that is, $V\subseteq U_v^\perp$ and $V\cap K_0=D_v$, whose kernel direction
is $\ker(\mathcal{L}|_{K_0})=\mathbb{R} v$, if and only if $v$ is an eigenvector of the compression:
$Mv=\mu v$ for some $\mu\in\mathbb{R}$. Consequently the set of directions satisfying
\ref{sec3:S1}--\ref{sec3:S2} is the union of the projective eigenspaces
$\mathbb{P}(E_\mu)$, $\mu\in\operatorname{Spec}_{\mathbb{R}}(M)$, and this structural admissibility of $v$
is equivalent to each of: $Mv\in\mathbb{R} v$; $\varphi_{Av}\in\mathbb{R}\varphi_v$; $Av-\mu v\in K_0^\perp$
for some real $\mu$. The theorem characterizes \ref{sec3:S1}--\ref{sec3:S2} only: \ref{sec3:S3}
($V_0\subseteq V$) is an extra condition. Claim~\ref{sec3:lem:prescribed} shows that it restricts the
admissible directions to $\mathbb{P}(E_\mu^0)$, and Claim~\ref{sec3:lem:constantrank} adds the fibre
condition $\delta\le s\le r(v)$; a direction satisfying \ref{sec3:S1}--\ref{sec3:S2} but lying outside
$E_\mu^0$ is therefore not a kernel direction occurring in the layer with the prescribed rows.
\end{thm}
\begin{proof}
\emph{(Sufficiency.)} If $Mv=\mu v$ then, for every $u\in K_0$,
\begin{align*}
\varphi_{Av}(u)=\langle u,Av\rangle
&=\langle u,\pi_{K_0}(Av)\rangle=\langle u,Mv\rangle\\
&=\mu\langle u,v\rangle=\mu\,\varphi_v(u),
\end{align*}
so $\varphi_{Av}=\mu\varphi_v$ and the two functionals have rank $1$ (as $v\neq0$ and the inner product
is positive definite, $\varphi_v\neq0$). By Claim~\ref{sec3:lem:direction}, $\dim D_v=k-1$. Choose $V$
with $D_v\subseteq V\subseteq U_v^\perp$ and $\dim V=q$, possible if and only if
$k-1\le q\le\dim U_v^\perp$, the nonemptiness condition of Claim~\ref{sec3:lem:constantrank}. Then
$\pi_{K_0}(V)\subseteq K_0\cap U_v^\perp=D_v$ and $\pi_{K_0}(V)\supseteq\pi_{K_0}(D_v)=D_v$, so
$\pi_{K_0}(V)=D_v$ and $\omega=k-1$, that is, $c=1$, while
$\ker(\mathcal{L}|_{K_0})=K_0\cap V^\perp=\mathbb{R} v$.

\emph{(Necessity.)} If $v$ is the kernel direction of a $V$ with \ref{sec3:S1}--\ref{sec3:S2}, then
$V\subseteq U_v^\perp$ and $\dim(V\cap K_0)=k-1$ by
\ref{sec3:S1}--\ref{sec3:S2} and Claim~\ref{sec3:lem:direction}(iv). Since
$V\cap K_0\subseteq K_0\cap U_v^\perp=D_v$, we get
$\dim D_v\ge k-1$; combined with Claim~\ref{sec3:lem:direction}(iii) this gives
$\rank\{\varphi_v,\varphi_{Av}\}\le1$. As
$\varphi_v\neq0$ there is $\mu\in\mathbb{R}$ with $\varphi_{Av}=\mu\varphi_v$, i.e.\
$\langle u,Av-\mu v\rangle=0$ for all $u\in K_0$, so $Av-\mu v\in K_0^\perp$. But
$Av-\mu v=(Mv-\mu v)+A_{21}v$ with $Mv-\mu v\in K_0$ and $A_{21}v\in K_0^\perp$, so its
$K_0$-component $Mv-\mu v$ must vanish: $Mv=\mu v$. Finally $\mu\in\mathbb{R}$ because $v$ and $M$ are
real.

Both directions concern the structural conditions only: the spectral condition \eqref{eq:5b} is an
independent restriction tested on the fibre, so an eigen-direction is structurally admissible but not
necessarily feasible.
\end{proof}


Every direction is therefore an eigenvector; what remains is to ask which eigenvectors the prescribed
rows leave.


\subsection{Directions compatible with the prescribed rows, and the fibre}\label{sec3:sub:prescribed}


The purpose of this step is to impose the prescribed rows once and for all, and to extract $\delta$
and $E_\mu^0$ from them; it is used in Claim~\ref{sec3:lem:constantrank}.

\begin{claim}[Prescribed rows]
\label{sec3:def:E0}
\label{sec3:lem:prescribed}
\label{sec3:cor:empty}
\label{sec3:proc}
\label{sec3:prop:rankcheck}
Let $\mu\in\operatorname{Spec}_{\mathbb{R}}(M)$ and $v\in E_\mu^0\setminus\{0\}$, where
$E_\mu^0=\{v\in E_\mu:\ L_0v=0,\ L_0Av=0\}$. Then:
\begin{enumerate}[label=(\roman*),leftmargin=*,itemsep=1pt,topsep=2pt]
\item $V_0\subseteq U_v^\perp$, $V_0\cap D_v=V_0\cap K_0$, and
$\overline{V}_{0,v}:=(V_0+D_v)/D_v\subseteq\mathcal{Q}_v$ has dimension
$\delta=\dim V_0-\dim(V_0\cap K_0)$, independent of $v$;
\item if $E_\mu^0=\{0\}$ for every $\mu\in\operatorname{Spec}_{\mathbb{R}}(M)$ then the layer is empty
and the instance is infeasible for every $q$;
\item if $\rank L_0<r_0$ then no $\mathcal{L}=[L_0;R]$ has full row rank, so the instance is
infeasible and the standing assumption $\rank L_0=r_0$ is forced, not imposed.
\end{enumerate}
\end{claim}
\begin{proof}
By Claim~\ref{sec3:lem:direction} a direction occurring in a design with prescribed rows must satisfy
$L_0v=L_0Av=0$, that is $v\in E_\mu^0$; conversely $L_0v=L_0Av=0$ says that every row of $V_0$
annihilates $v$ and $Av$, i.e.\ $V_0\subseteq U_v^\perp$. Hence
$V_0\cap D_v=V_0\cap K_0\cap U_v^\perp=V_0\cap K_0$, which does not involve $v$, so $\delta$ is fixed
and $\dim\overline{V}_{0,v}=\dim V_0-\dim(V_0\cap D_v)=\delta$. For (ii), the kernel direction of a
design in the layer belongs to $E_\mu^0$ for some real $\mu$ by Claim~\ref{sec3:lem:direction}; no
direction, no design. For (iii), $\dim V\le\rank L_0+(q-r_0)<q$ contradicts $\dim V=q$.
\end{proof}


With the prescribed rows imposed, a direction carries a dichotomy that changes the size of the fibre
above it.

\begin{prop}[Eigen and non-eigen directions]
\label{sec3:lem:eigen}
\label{sec3:prop:strata}
\label{sec3:lem:dimU}
\label{sec3:rem:rankcases}
Let $\mu\in\operatorname{Spec}_{\mathbb{R}}(M)$ and $v\in E_\mu\setminus\{0\}$. Then:
\begin{enumerate}[label=(\roman*),leftmargin=*,itemsep=1pt,topsep=2pt]
\item $Av\in\operatorname{span}(v)\iff A_{21}v=0$; the \emph{eigen directions} form the subspace
$E_\mu^{\mathrm{eig}}=E_\mu\cap\ker A_{21}$, whose projectivization
$\mathbb{P}(E_\mu^{\mathrm{eig}})\subseteq\mathbb{P}(E_\mu)$ is a projective subspace, while the
\emph{non-eigen directions} form its complement, a nonempty Zariski-open subset whenever
$E_\mu^{\mathrm{eig}}\ne E_\mu$;
\item if $A_{21}v=0$ then $U_v=\operatorname{span}(v)$ and $\dim U_v=1$, whereas if $A_{21}v\ne0$ then
$v,Av$ are independent and $\dim U_v=2$; accordingly $\dim U_v^\perp=n-1$ or $n-2$;
\item with prescribed rows the two branches are $B_\mu^{\mathrm{eig}}=\mathbb{P}(E_\mu^0\cap\ker A_{21})$
and $B_\mu^{\mathrm{non}}=\mathbb{P}(E_\mu^0)\setminus B_\mu^{\mathrm{eig}}$, which are disjoint and
cover $\mathbb{P}(E_\mu^0)$; according to $\rank(A_{21}|_{E_\mu^0})$ the second is empty (rank zero),
the first is empty (full rank), or both are nonempty—the first screening step of
Section~\ref{sec4}.
\end{enumerate}
\end{prop}
\begin{proof}
$Av=Mv+A_{21}v=\mu v+A_{21}v$ with $\mu v\in K_0$ and $A_{21}v\in K_0^\perp$, so
$Av\in\operatorname{span}(v)\subseteq K_0$ if and only if the $K_0^\perp$-component $A_{21}v$ vanishes.
As $\ker A_{21}$ is linear, its intersection with $E_\mu$ is linear and projectivization preserves the
linear structure; the complement is open and nonempty if and only if
$E_\mu^{\mathrm{eig}}\ne E_\mu$. For (ii), $A_{21}v=0$ gives $Av=\mu v$; otherwise $Av-\mu v\ne0$ and
$\langle v,Av-\mu v\rangle=\langle v,A_{21}v\rangle=0$ (as $A_{21}v\in K_0^\perp$ and $v\in K_0$), so
$v,Av$ are independent; the orthogonal complement has dimension $n-\dim U_v$. The two branches differ
in $\dim U_v$, hence in the rank of the quotient bundle and in the fibre dimension, which is the
content of Claim~\ref{sec3:lem:constantrank}.
\end{proof}


The purpose of this step is the fibre: it identifies the row spaces over a fixed direction with a
Grassmannian and computes its dimension; it is used in Theorem~\ref{sec3:thm:param}.

\begin{claim}[Constant rank and the fibre]
\label{sec3:lem:constantrank}
\label{sec3:lem:fibre}
\label{sec3:lem:fullrank}
Over $\mathbb{P}(E_\mu^0)$ the assignment $v\mapsto D_v$ is a vector bundle of constant rank $k-1$. The
quotient $\mathcal{Q}_v=U_v^\perp/D_v$ is \emph{not} of constant rank over all of $\mathbb{P}(E_\mu^0)$:
its rank is $\rho$ on the eigen stratum $\mathbb{P}(E_\mu^0\cap\ker A_{21})$ and $\rho-1$ on the
non-eigen stratum, and it is a vector bundle of that constant rank on each stratum separately. For a
fixed direction
$[v]$, the row spaces satisfying \ref{sec3:S1}--\ref{sec3:S3} with $\dim V=q$ are exactly
$\{V:\ V_0+D_v\subseteq V\subseteq U_v^\perp,\ \dim V=q\}$, and $V\mapsto V/D_v$ identifies this set
with $\Gr\bigl(s-\delta,\ \mathcal{Q}_v/\overline{V}_{0,v}\bigr)$, of dimension
$(s-\delta)\bigl(r(v)-s\bigr)$, nonempty if and only if $\delta\le s\le r(v)$. Moreover, if a design
in the layer exists with $\dim V\le q$ and $q\le\dim U_v^\perp$, then a design exists with $\dim V=q$
and the same compressed pencil conditions.
\end{claim}
\begin{proof}
By Theorem~\ref{sec3:thm:direction} the two functionals are proportional on $\mathbb{P}(E_\mu^0)$, so
their rank is constantly $1$ and $\dim D_v$ is constantly $k-1$. Moreover
$\dim\mathcal{Q}_v=\dim U_v^\perp-(k-1)=(n-\dim U_v)-(k-1)$, and $\dim U_v=1$ or $2$ according to the
two strata of Proposition~\ref{sec3:prop:strata}; substituting $\rho=n-k$ gives the stated values, and
constant rank gives the bundle structure. In particular $r(v)\le\rho\le2p$, the starting point of
Section~\ref{sec5}.

If $D_v\subseteq V\subseteq U_v^\perp$ then $V\cap K_0\subseteq K_0\cap U_v^\perp=D_v$, with the
reverse inclusion from $D_v\subseteq V$, so \ref{sec3:S2} holds automatically; conversely
\ref{sec3:S2} with $V\subseteq U_v^\perp$ gives $V\cap K_0=D_v$. Hence $V/D_v$ is an $s$-dimensional
subspace of $\mathcal{Q}_v$ containing $\overline{V}_{0,v}=(V_0+D_v)/D_v$. The $s$-dimensional
subspaces containing a fixed $\delta$-dimensional one correspond to the $(s-\delta)$-dimensional
subspaces of $\mathcal{Q}_v/\overline{V}_{0,v}$, which is the Grassmannian above; its dimension is the
standard formula $(s-\delta)(r(v)-s)$ and its nonemptiness condition is $0\le s-\delta\le r(v)-\delta$.
The map $V\mapsto V/D_v$ is injective with exactly that image, so no admissible row space is lost.

For the normalization, choose $\tilde V$ with $V\subseteq\tilde V\subseteq U_v^\perp$ and
$\dim\tilde V=q$, possible because $q\le\dim U_v^\perp$. Then \ref{sec3:S1}--\ref{sec3:S3} and
$\tilde V\cap K_0=D_v$ are preserved, so $\rank\tilde H=k-1=\omega$ and
$\row\tilde H_A\subseteq\row\tilde H$; and \eqref{eq:5b} is preserved because enlarging the row space
only adds rows to $P(\lambda)$, which cannot lower the rank below the structural bound $\omega$. We may
therefore work with $q$-dimensional row spaces throughout: the row count and $\rank L_0=r_0$ are the
only constraints on the dimension, and the normalization does not change the compressed pencil.
\end{proof}


With the rows imposed and the fibre computed, the pieces can be assembled into one parameter space.


\subsection{The complete parameterization}\label{sec3:sub:param}


The purpose of this step is to assemble the strata and their fibres into the object on which the
decision procedure runs.

\begin{thm}[Complete parameterization]
\label{sec3:thm:param}
\label{sec3:rem:empty}
\label{sec3:rem:dimension}
\label{sec3:rem:realpoints}
Assume $\rank L_0=r_0$ and $c=1$. Then the parameter space of row spaces of the hyperplane layer is
the finite union
\begin{gather*}
\bigsqcup_{\mu\in\operatorname{Spec}_{\mathbb{R}}(M)}
\Bigl(
\Gr_{B_\mu^{\mathrm{eig}}}\bigl(s-\delta,\ \mathcal{Q}^{\mathrm{eig}}/\overline{V}_0\bigr)
\\
\sqcup\ \Gr_{B_\mu^{\mathrm{non}}}\bigl(s-\delta,\ \mathcal{Q}^{\mathrm{non}}/\overline{V}_0\bigr)
\Bigr),
\end{gather*}
where $B_\mu^{\mathrm{eig}},B_\mu^{\mathrm{non}}$ are as in Proposition~\ref{sec3:prop:strata}, the
bundles $\mathcal{Q}^{\mathrm{eig}},\mathcal{Q}^{\mathrm{non}}$ are the constant-rank quotient bundles
of Claim~\ref{sec3:lem:constantrank}, and $\overline{V}_0$ is the subbundle of
Claim~\ref{sec3:lem:prescribed}. Each stratum is an \emph{algebraic relative Grassmann bundle over a
locally closed real algebraic base}, of dimension
$\dim X_\mu^{\mathrm{eig}}=h_\mu^0-1+(s-\delta)(\rho-s)$ and
$\dim X_\mu^{\mathrm{non}}=d_\mu^0-1+(s-\delta)(\rho-1-s)$, with nonemptiness conditions
$\delta\le s\le\rho$ on the eigen stratum and $\delta\le s\le\rho-1$ on the non-eigen stratum. Every
nonempty stratum is irreducible and its real points are Zariski dense. The dimension displayed above is
the sum of a base term and a fibre term: the fibre term is
$(s-\delta)(r(v)-s)\le r(v)^2/4\le p^2$, and the base term is $d_\mu^0-1$ (or $h_\mu^0-1$), which
vanishes exactly in the geometrically simple case $d_\mu^0=1$
(respectively $h_\mu^0=1$) of Section~\ref{sec5}. The total is therefore $\le p^2$ under that promise;
without it the base term grows with $n$ and no bound in $p$ alone is claimed. Emptiness is decided by
$(\delta,s,\rho,h_\mu^0,d_\mu^0)$ alone: if $h_\mu^0=0$ the eigen
stratum is empty, and if $d_\mu^0=0$ both are.
\end{thm}
\begin{proof}
By Theorem~\ref{sec3:thm:direction} the kernel direction lies in $\mathbb{P}(E_\mu)$ and, with
prescribed rows, in $\mathbb{P}(E_\mu^0)$ by Claim~\ref{sec3:lem:prescribed};
Proposition~\ref{sec3:prop:strata} splits each $\mathbb{P}(E_\mu^0)$ into the two strata;
Claim~\ref{sec3:lem:constantrank} gives the fibre over each stratum and shows its dimension constant on
the stratum, so the fibres form a relative Grassmann bundle. The bases are a projective space or an open
subset of one, both irreducible over the reals, and a bundle with irreducible base and constant fibre
dimension is irreducible; the dimension is that of the base plus that of the fibre, with base
dimensions $h_\mu^0-1$ and $d_\mu^0-1$. The nonemptiness conditions follow by substituting the two
values of $r(v)$ into $\delta\le s\le r(v)$. The bases and fibres carry rational real charts whose real
points are Zariski dense, which is what lets Section~\ref{sec4} pass from a function-field condition to
a real design.
\end{proof}


\paragraph*{Duality with the compressed formulation}\label{sec3:sub:duality}

Let $\mathsf{M}=[N\ \ AN]$ and let $\Phi\colon\mathbb{R}^n\to\mathbb{R}^{k}\oplus\mathbb{R}^{k}$,
$x\mapsto(xN,xAN)$, in the row-vector convention; for a row space $V$ write
$W=\Phi(V)=\row(V\mathsf{M})\subseteq\mathbb{R}^{2k}$, let $\pi_1,\pi_2$ be the coordinate projections,
and for $a\in\mathbb{R}^k$ put $a^\circ=\{h\in(\mathbb{R}^k)^*: h(a)=0\}$. In \cite{companion2} a
feasible design is described by $W$ through $\pi_2(W)\subseteq\pi_1(W)$ and
$\dim\pi_\lambda(W)=\dim\pi_1(W)$ for all $\lambda$, with $\pi_1(W)=\pi_2(W)=U$, $\dim U=\omega$ and
$\omega+1\le\dim W\le 2\omega$.

\begin{prop}[Direction duality and the fibres of $V\mapsto W$]
\label{sec3:prop:duality}
\label{sec3:prop:fibres}
\label{sec3:rem:noalternative}
Let $V$ be a row space of the layer with direction $v=Na$. Then $W\subseteq a^\circ\times a^\circ$; if
moreover $\pi_1(W)=\pi_2(W)=U$ with $\dim U=k-1$, then $U=a^\circ$ and
$[a]=\mathbb{P}(\Ann(U))$, so the common projection of the compressed formulation and the direction of
this section are dual descriptions of the same object. As for the fibre of the map $V\mapsto W$:
$\ker\Phi=\{x:\ xN=0,\ xAN=0\}=(K_0+AK_0)^\perp$, and with $T_W=\Phi^{-1}(W)$, $S=V_0+D_v\subseteq T_W$,
the map $\Phi$ induces $\overline{\Phi}\colon T_W/S\to W/\Phi(S)$, the lifts of $W$ being exactly the
subspaces $\overline{V}=V/S\subseteq T_W/S$ with $\overline{\Phi}(\overline{V})=W/\Phi(S)$ and
$\dim\overline{V}=q-\dim S$; the set of lifts is thus a \emph{surjectivity locus} inside a
Grassmannian, in general not the whole Grassmannian. Since $\ker\Phi$ is in general nonzero, the
compressed formulation is more economical for deciding feasibility, while the ambient one retains what
the prescribed rows and the row count need.
\end{prop}
\begin{proof}
Since $V\subseteq U_v^\perp$, every row $x\in V$ satisfies $xv=0$ and $xAv=0$; with $v=Na$ this reads
$xNa=0$ and $xANa=0$, so $W\subseteq a^\circ\times a^\circ$ and hence
$\pi_1(W),\pi_2(W)\subseteq a^\circ$. If both equal $U$ with $\dim U=k-1=\dim a^\circ$, then
$U=a^\circ$; the annihilator of $a^\circ$ is one-dimensional in $(\mathbb{R}^k)^{**}\cong\mathbb{R}^k$
and spanned by $a$, so the projective point $[a]$ is determined by $U$.

The description of $\ker\Phi$ is immediate from the definition, and $xN=xAN=0$ is equivalent to
$x\perp K_0$ and $x\perp AK_0$. If $V\subseteq T_W$ then $\Phi(V)=W$ if and only if the image of $V$
under $\Phi$ equals $W$; after passing to the quotient by $S$ this becomes
$\overline\Phi(\overline V)=W/\Phi(S)$, and the dimension condition follows from $\dim V=q$ and
$\dim S$.
\end{proof}


The parameter space is therefore an explicit finite union of bundles, and the next section decides
feasibility on it.


\section{The decision procedure}\label{sec4}\label{sec:decision}

This section turns the parameterization into an exact decision procedure: finitely many charts, one
gcd computation per chart, and a leading-coefficient claim that carries the verdict at the generic
point to every point of the chart.

The compressed criterion of Section~\ref{sec:prelim} is decided by a greatest-common-divisor test on
the maximal minors. The closedness of the bad locus, the generic-point criterion and the finite-cover
principle that accompany it are standard facts about a family of polynomials over an irreducible base,
recalled below and \emph{not} taken from \cite{companion1}: the closedness is proved in
Appendix~\ref{appBC} and the step from a generic point to the whole chart is Claim~\ref{sec4:lem:fullrank}
here. What they do not supply on their own, and what this section develops, is the passage from a
\emph{generic} point to an \emph{entire chart}, so that one computation per chart gives a verdict valid at
every point of it. In the hyperplane layer that passage is supplied by a single invariant,
$\rank H\equiv\omega$, which keeps a common factor from degenerating to a constant
(Claim~\ref{sec4:lem:fullrank}) --- the counterpart, in the chart, of the use of the same invariant to keep
common roots from escaping to infinity. No complexity assertion is made
here; that is the content of Section~\ref{sec5}.


\subsection{Charts, and the leading-coefficient claim}\label{sec4:charts}


The purpose of this step is to fix what a chart is and what a parameter point on it looks like.

\begin{defn}[Good points, and affine charts of the fibre]
\label{sec4:defn:goodbad}
\label{sec4:defn:charts}
A parameter point $X$ is \emph{good} if $\rank P(X,\lambda)=\omega$ for every $\lambda\in\mathbb C$,
and \emph{bad} otherwise; by the criterion of Section~\ref{sec:prelim} it is bad exactly when
$\deg_\lambda\gcd_I m_I(X,\cdot)>0$. The charts are constructed as follows. Fix an admissible direction
$[v]\in\mathbb P(E_\mu^0)$ and a basis $e_1,\dots,e_{\widetilde r}$ of
$\mathcal Q_v/\overline V_{0,v}$. For $I\subseteq\{1,\dots,\widetilde r\}$ with $|I|=m=s-\delta$ put
$J=\{1,\dots,\widetilde r\}\setminus I$ and let
$\mathcal X_I\cong\mathbb C^{m\times(\widetilde r-m)}$ be the affine chart with coordinates
$X=(X_{ij})_{i\in I,\,j\in J}$. The row space attached to $X$ is spanned by the base rows --- a basis
of $V_0+D_v$, of which there are $(k-1)+\delta$ --- together with
\[
\row_i=e_i+\sum_{j\in J}X_{ij}\,e_j,\qquad i\in I .
\]
These are $(k-1)+\delta+m=(k-1)+s=q$ rows and they form $L(X)$. The chart count is
$\binom{\widetilde r}{m}$, the parameter count is
$d=m(\widetilde r-m)=(s-\delta)(r(v)-s)\le p^2$, and for $m=0$ the chart is a single point.
\end{defn}

The charts above are \emph{affine} and \emph{explicitly} parameterized: the rows are affine in $X$, so
the minors are polynomials a computer can form, and each chart is irreducible by construction, so
Claim~\ref{sec4:thm:generic} applies without any decomposition. What they do not yet provide is a
cover of the \emph{total} parameter space: when $\dim E_\mu^0>1$ the direction is itself a parameter and
the rows of $L$ depend on it through $D_v$ and through a basis of the quotient. The next definition
supplies the missing charts, so that the algorithm of Section~\ref{sec4:algorithm} enumerates a finite
list rather than the points of a continuum.

\begin{defn}[Charts of the total parameter space]
\label{sec4:defn:totcharts}
Fix a real eigenvalue $\mu$ with $K=\mathbb Q(\mu)$ and the real embedding selecting $\mu$, let
$d_\mu^0=\dim E_\mu^0$, $h_\mu^0=\dim(E_\mu^0\cap\ker A_{21})$, and fix a $K$-basis
$e_1,\dots,e_{d_\mu^0}$ of $E_\mu^0$. A chart of the total parameter space is a product of three
factors, each a principal open subset of an affine space over $K$.
\begin{enumerate}[label=(\roman*),leftmargin=*]
\item \emph{Direction chart.} If $d_\mu^0=1$ the factor is a point. Otherwise, for each
$j\in\{1,\dots,d_\mu^0\}$ let $v_j(x)=e_j+\sum_{i\ne j}x_i e_i$ with coordinates $x\in K^{d_\mu^0-1}$;
these $d_\mu^0$ charts cover $\mathbb P(E_\mu^0)$. The eigen stratum
$\mathbb P(E_\mu^0\cap\ker A_{21})$ is a projective subspace of dimension $h_\mu^0-1$; it is covered by
the $h_\mu^0$ charts obtained in the same way from a $K$-basis of $E_\mu^0\cap\ker A_{21}$. The non-eigen
stratum is covered by the intersections of the $v_j$-charts with the opens
$\{(A_{21}v_j(x))_q\ne0\}$, one for each row index $q$ of $A_{21}$; the factors
$\{(A_{21}v_j(x))_q\}$ are the pivot conditions of the trivialization.
\item \emph{Trivialization chart.} On each direction chart the two functionals $\varphi_v,\varphi_{Av}$
have entries that are affine in $x$, so $D_v=K_0\cap U_v^\perp$, $U_v^\perp$ and the quotient
$\mathcal Q_v/\overline V_{0,v}$ are obtained by Gaussian elimination over $K(x)$; the pivots are
specified minors of the $2\times k$ matrix $[\varphi_v;\varphi_{Av}]$ and of the analogous
$2\times n$ matrix, and each choice of pivots gives a basis whose entries are rational functions of $x$
with denominators equal to those minors. Refining the direction chart by those nonvanishing conditions
gives a locally closed affine chart on which the basis is polynomial after clearing denominators.
\item \emph{Fibre chart and rank chart.} On the trivialized chart, add a Grassmann chart of
Definition~\ref{sec4:defn:charts} for the fibre, and refine by the condition that a specified
$\omega\times\omega$ minor $\Delta$ of $H$ is nonzero; on that open set $\rank H=\omega$ identically.
\end{enumerate}
\end{defn}


The purpose of this step is to make the chart list finite and exhaustive; it is used in the
correctness proof of Algorithm~\ref{sec4:algo:decide}.

\begin{claim}[The total charts cover, and are finitely many]\label{sec4:lem:cover}
Every admissible point of the layer lies in at least one chart of Definition~\ref{sec4:defn:totcharts};
the number of charts is finite, and on each of them the entries of $L$, $H$ and $H_A$ are rational
functions of the joint coordinates whose denominators are the pivot and rank minors, while
$\rank H\equiv\omega$ holds identically.
\end{claim}

\begin{proof}
An admissible point determines a real eigenvalue $\mu$, a direction $[v]\in\mathbb P(E_\mu^0)$, a point
of the fibre over $[v]$ and a positive value of $\rank H$ equal to $\omega$; the direction lies in one
of the direction charts, the elimination pivots are those of its own trivialization, some
$\omega$-minor $\Delta$ of $H$ is nonzero there, and the fibre point lies in one Grassmann chart. Hence
the point lies in the corresponding product chart. Finiteness follows because each of the three factors
is indexed by a finite set: $d_\mu^0$ (or $h_\mu^0$) projective charts, finitely many pivot choices,
finitely many $\omega$-minors and $\binom{\widetilde r}{m}$ Grassmann charts, and there are at most $k$
real eigenvalues. The rationality statement is the content of (ii) and (iii): the elimination is
performed over $K(x)$, so the basis entries are ratios of minors, and clearing denominators makes them
polynomial on the refined open set; $\rank H=\omega$ is imposed by $\Delta\ne0$.
\end{proof}


The purpose of this step is to bound the size of the objects one chart carries; it is used in
Section~\ref{sec5}.

\begin{claim}[Size of a chart]
\label{sec4:lem:minorcount}
On every chart of Definition~\ref{sec4:defn:charts} the entries of $L(X)$, hence of $H(X)$ and
$H_A(X)$, are affine in $X$; every maximal minor satisfies $\deg_X m_I\le\omega$ and
$\deg_\lambda m_I\le\omega$; and the number of maximal minors is
$\binom{q}{\omega}\binom{k}{\omega}=\binom{q}{k-1}k=k\binom{q}{s}$.
\end{claim}

\begin{proof}
$H(X)=L(X)N$ is affine in $X$ because $N$ is constant, and $m_I$ is an $\omega\times\omega$
determinant, hence multilinear, so its degree is at most $\omega$ in each of $X$ and $\lambda$. The
count is the number of ways of choosing $\omega$ of the $q$ rows and $\omega$ of the $k$ columns,
with $\binom{k}{k-1}=k$ and $\binom{q}{k-1}=\binom{q}{s}$.
\end{proof}

\emph{Why a generic test is not yet a decision.} Claim~\ref{sec4:thm:generic} tests the generic point
of an irreducible piece, which answers a question about almost every point, whereas a decision
procedure needs an answer at every point: a chart whose generic point is bad may in principle contain
good points. The gap is real for general families; what closes it here is the following, the main new
result of this section.

The coefficient of $\lambda^\omega$ in $m_I(X,\lambda)$ is the $\omega\times\omega$ minor of $H(X)$
with the same index set. Since $\rank H(X)=\omega$ at every point of a chart, at least one of these
minors is nonzero, that is,
\begin{equation}
\label{sec4:eq:fullrank}
\text{for every }X\in\mathcal X_I\ \text{there is }I\text{ with }
\deg_\lambda m_I(X,\cdot)=\omega .
\end{equation}


The purpose of this step is to pass from the generic point of a chart to every point of it; it is used
wherever a function-field verdict is read as a verdict on designs.

\begin{claim}[Full-rank leading coefficient; dichotomy on a chart]
\label{sec4:lem:fullrank}
\label{sec4:cor:dichotomy}
Let $K$ be a field, $X=(X_1,\dots,X_d)$, and let $m_1,\dots,m_t\in K[X,\lambda]$ satisfy
$\deg_\lambda m_i\le\omega$ and \eqref{sec4:eq:fullrank}. If
$\deg_\lambda\gcd(m_1,\dots,m_t)>0$ in $K(X)[\lambda]$, then for \emph{every} $X_0$ the specialized
polynomials $m_i(X_0,\cdot)$ have a nonconstant common factor. In particular, on a chart of the layer
exactly one of the following holds: either $\deg_\lambda\gcd_I m_I=0$, in which case the chart contains
a good point and its good points form a nonempty Zariski open subset; or
$\deg_\lambda\gcd_I m_I>0$, in which case \emph{every} point of the chart is bad.
\end{claim}

\begin{proof}
Let $g$ be the gcd in $K(X)[\lambda]$, of positive degree. By Gauss's lemma \cite{lan02} we may take
$g$ primitive in $K[X][\lambda]$ and of positive degree, with $g\mid m_i$ in $K[X][\lambda]$ for every
$i$; write $m_i=g\,h_i$ and $\omega_i=\deg_\lambda m_i$, $\gamma=\deg_\lambda g\ge1$. Fix $X_0$. If
$g(X_0,\cdot)\equiv0$ then $m_i(X_0,\cdot)=g(X_0,\cdot)h_i(X_0,\cdot)\equiv0$ for every $i$,
contradicting \eqref{sec4:eq:fullrank}. If $g(X_0,\cdot)$ is a nonzero constant then
$\deg_\lambda m_i(X_0,\cdot)\le\deg_\lambda h_i\le\omega_i-\gamma\le\omega-\gamma<\omega$ for every
$i$, again contradicting \eqref{sec4:eq:fullrank}. Hence $g(X_0,\cdot)$ is nonzero and nonconstant, and
it is a common factor of all $m_i(X_0,\cdot)$. The dichotomy follows by applying this to the charts and
Claim~\ref{sec4:thm:generic} to the degree-zero alternative.
\end{proof}

Without \eqref{sec4:eq:fullrank} the implication is false: $m(X,\lambda)=X\lambda+1$ has gcd $m$, of
degree $1>0$ over $K(X)[\lambda]$, yet at $X_0=0$ it specializes to the constant $1$, and
\eqref{sec4:eq:fullrank} fails there. The lemma thus rests on the invariant $\rank H\equiv\omega$,
which is what the layer supplies.


Charts and this claim give a verdict chart by chart; the next subsection turns them into certificates
and into an algorithm.


\subsection{Certificates and the decision algorithm}\label{sec4:real}\label{sec4:certificates}\label{sec4:algorithm}


The purpose of this step is to say what a returned certificate consists of and which alternative
establishes which verdict; it is used in the correctness proof of the algorithm.

\begin{claim}[Real designs and certificates]
\label{sec4:prop:realpoints}
\label{sec4:prop:sieve}
\label{sec4:prop:dimineq}
\label{sec4:prop:resultant}
\label{sec4:prop:cert}
The following hold, in decreasing logical strength.
\begin{enumerate}[label=(\alph*),leftmargin=*]
\item \label{it:real} Every chart $\mathcal X_I\cong\mathbb C^d$ carries the standard real structure
$\mathbb R^d$, and so do the affine charts of $\mathbb P(E_\mu^0)$ because $\mu$ is real and $E_\mu^0$
is defined by real linear equations. The real points of every chart are Zariski dense in it; if the
good set of a chart is a nonempty Zariski open subset then it contains a real point, hence a real
admissible design; and if $\deg_\lambda\gcd>0$ on a chart then every real point of it is bad.
\item \label{it:sieve} If $E_\mu^0=\{0\}$ for every real eigenvalue $\mu$ of
$M=\pi_{K_0}A|_{K_0}$, then the layer is infeasible.
\item \label{it:dimineq} If the fibre over some $[v]\in\mathbb P(E_\mu^0)$ is nonempty then
$\delta\le s\le r(v)$; if $\delta>s$, or $s>\rho$ on the eigen branch and $s>\rho-1$ on the non-eigen
branch, the layer is infeasible.
\item \label{it:resultant} If $\operatorname{Res}_\lambda(m_I,m_J)\not\equiv0$ in $K[X]$ for some pair
$I,J$, then $\deg_\lambda\gcd_I m_I=0$ and the chart, hence the layer, is feasible. The converse fails:
$f_1=\lambda(\lambda-1)$, $f_2=\lambda(\lambda-2)$ and $f_3=(\lambda-1)(\lambda-2)$ share a factor
pairwise, respectively $\lambda$, $\lambda-1$ and $\lambda-2$, yet $\gcd(f_1,f_2,f_3)=1$; a nonzero
resultant is a usable positive certificate, whereas failure to find one is not a negative certificate.
\item \label{it:cert} A complete decision certificate consists of (i) the factorization of $\chi_M$ and
a basis of $E_\mu^0$ over each irreducible factor, or the statement $E_\mu^0=0$; (ii) for every
admissible direction and chart, the chart and parameter counts; (iii) for every chart either a
\emph{good} alternative --- a Bézout identity $\sum_I a_I m_I=1$ with $a_I\in K[X][\lambda]$, checked
by one multiplication --- or a \emph{bad} alternative --- a nonconstant $g\in K[X][\lambda]$ with
division identities $g\,h_I=m_I$ for every $I$. If some chart reports ``good'' the layer is feasible;
if every chart reports ``bad'' it is infeasible.
\end{enumerate}
\end{claim}

\begin{proof}
(a) The zero set of a nonzero polynomial has empty interior in $\mathbb R^d$, so its complement is
dense \cite{bpr06}; the good set is a non-vanishing set of a nonzero polynomial, so it meets
$\mathbb R^d$, and the corresponding $L(X)$ is a real matrix with a real row space.
(b) Every admissible direction lies in some $E_\mu^0$.
(c) The fibre is $\Gr(m,\widetilde r)$ with $m=s-\delta$, nonempty exactly when
$0\le m\le\widetilde r$.
(d) A nonzero resultant makes $m_I$ and $m_J$ coprime in $K(X)[\lambda]$, so the gcd is $1$; the
failure of the converse is the example displayed.
(e) The good alternative certifies $\gcd=1$ without any gcd computation; for the bad alternative,
Claim~\ref{sec4:lem:fullrank} shows that $g$ is nonconstant and divides every minor at every point, so
the whole chart is bad, and the cover is finite.
\end{proof}


The certificate turns the criterion into a test that can be checked once per chart; what remains is to
enumerate the charts and apply it.

\begin{algo}[Exact decision for the entire layer]
\label{sec4:algo:decide}
\begin{enumerate}
\item \emph{Pre-check.} If $\rank L_0<r_0$, output ``infeasible'' and stop: $\mathcal L$ must have
full row rank, so $\rank L_0=r_0$ is necessary.
\item Compute $K_0$, $N$ and $M=\pi_{K_0}A|_{K_0}$; factor $\chi_M$ and keep its real irreducible
factors.
\item For every real irreducible factor, enumerate the charts of Definition~\ref{sec4:defn:totcharts}:
first the direction charts of \ref{sec4:defn:totcharts}(i), taking the eigen charts when
$h_\mu^0\ge1$ and the non-eigen charts when $h_\mu^0<d_\mu^0$, and compute $\delta,s,m,\widetilde r$ on
each; if $\delta>s$ or $s>r(v)$, skip that chart by Claim~\ref{sec4:prop:dimineq}. This
enumeration is finite by Claim~\ref{sec4:lem:cover} and does \emph{not} sample individual directions.
\item For every chart produced in step 3: trivialize the quotient as in
Definition~\ref{sec4:defn:totcharts}(ii), attach a fibre chart \ref{sec4:defn:totcharts}(iii), build
$L(X)$ over $K(x)$ with $x$ the joint coordinates, clear the denominators of the pivot and rank minors,
form the $k\binom{q}{s}$ maximal minors of $P(X,\lambda)$, and compute their gcd over $K(x)[\lambda]$
by the subresultant algorithm \cite{vzgg03}.
\item If some chart has $\deg_\lambda\gcd=0$, output ``feasible'' and stop.
\item If every chart has $\deg_\lambda\gcd>0$, output ``infeasible''.
\item If the computation cannot be completed, output ``undecided'', never ``infeasible''.
\end{enumerate}
\end{algo}

\begin{thm}[Correctness]
\label{sec4:prop:correct}
Whenever steps 1--6 are completed, the verdict of Algorithm~\ref{sec4:algo:decide} is correct.
\end{thm}

\begin{proof}
Step 1 is the problem definition and step 3 is Claim~\ref{sec4:prop:dimineq}. Step 3 is finite by
Claim~\ref{sec4:lem:cover}, and its charts cover the whole parameter space, so no admissible design is
lost. In step 5 a chart with $\deg=0$ contains a good point by Claim~\ref{sec4:thm:generic}; in step 6
Claim~\ref{sec4:cor:dichotomy} makes every chart entirely bad, and the cover is finite, so no good
point exists. The ``undecided'' output of step 7 is not a verdict. Two implementation points belong to
the proof as well: step 2 enumerates irreducible factors rather than roots, since eigenspaces of
conjugate roots are Galois conjugate and conjugation preserves the $\lambda$-degree of the gcd; and
all linear algebra in step 4 must be done over the \emph{same} field $K=\mathbb Q(\mu)$, because when
$\mu\notin\mathbb Q$ the bases of $E_\mu^0$, $D_v$, $U_v^\perp$ and
$\mathcal Q_v/\overline V_{0,v}$ have no rational coordinates, and one coefficient field is what makes
the chain ``chart rows, minors, function-field gcd'' correct.
\end{proof}


The procedure is exact and complete for the layer; what it does not say is how long it takes.

\section{Fixed-\texorpdfstring{$p$}{p} bit complexity}\label{sec5}\label{sec:complexity}


This section prices the procedure: a bit model for the number field, the size bounds it needs, and the
fixed-$p$ theorem. The exponent depends on $p$ only, and the promise is what reduces the direction
parameter from a projective eigenspace to a finite list of points; the procedure itself is the exact one
of Section~\ref{sec4} and is defined on every input.

This section turns the geometry and the decision procedure into a \emph{bit complexity} statement: for
fixed output dimension $p$ and a \emph{geometrically simple real spectrum}, deciding feasibility in the
hyperplane layer takes time polynomial in the bit size of the input. This is the only complexity result
of the paper and the only one carrying a promise; its boundaries are examined in the closing paragraph of this
section. Throughout, $\tau$ is the encoding length of $(A,C,L_0,q)$, all
arithmetic is exact, and proofs of the estimates are in Appendix~\ref{appBC}, where the explicit
bounds are recorded.


\subsection{The model and the bounds}\label{sec5:sub:gsimple}

\begin{defn}[Geometrically simple real spectrum]\label{sec5:def:gsimple}
Let $M=\pi_{K_0}A|_{K_0}$ be the compression. We say that $M$ has a \emph{geometrically simple real
spectrum} if $\dim\ker(M-\mu I)=1$ for every real eigenvalue $\mu$ of $M$.
\end{defn}

Each real eigenvalue then contributes exactly \emph{one} projective direction. The condition is weaker
than algebraic simplicity: an eigenvalue may have algebraic multiplicity $>1$ while its eigenspace
stays one-dimensional, and it still contributes a single direction, the base
$\mathbb{P}(E_\mu)$ being a point. Conversely, if some real eigenvalue has
$d_\mu=\dim E_\mu\ge2$, the base has dimension $d_\mu-1$ and the number of parameters grows with $n$
(the closing paragraph of this section). The prescribed rows enter through
$E_\mu^0=E_\mu\cap\ker(L_0N)\cap\ker(L_0AN)$: the enumerable directions are those of $E_\mu^0$, so
under geometric simplicity each real eigenvalue yields at most one candidate, and an eigenvalue with
$E_\mu^0=\{0\}$ is excluded outright.

All exact computations of the decision procedure take place in the number field generated by an
eigenvalue, and the bit bounds below refer to a fixed representation of that field. The representation
is chosen once and for all.

\begin{defn}[Effective representation of $K$]\label{sec5:def:model}
Let $\mu$ be a real eigenvalue of $M$ and let $f_K\in\mathbb{Z}[\theta]$ be the minimal polynomial of
$\mu$ over $\mathbb{Q}$, taken primitive with positive leading coefficient. Put
$d_K=\deg f_K\le k\le n$ and let $\beta_K$ be the largest bit length of the coefficients of $f_K$. The
field is $K=\mathbb{Q}(\mu)\cong\mathbb{Q}[\theta]/(f_K)$, and an element of $K$ is represented by its
coordinate vector in the power basis $(1,\theta,\dots,\theta^{d_K-1})$, each coordinate a rational.
The \emph{size} of an element is the largest bit length of its coordinates and is written $\beta$.
\end{defn}

Three points make this representation usable in a bit model.

\emph{(i) Field parameters.} By Claim~\ref{sec5:lem:height}(c) below, $d_K\le n$ and
$\beta_K=\poly(n,\tau)$: $f_K$ is a factor of the integer polynomial $\chi_M$, and the height of a
factor is bounded by the height and the degree of the polynomial it divides.

\emph{(ii) Cost of arithmetic.} Addition and multiplication are polynomial arithmetic in
$\mathbb{Q}[\theta]$ followed by reduction modulo $f_K$: $O(d_K^{2})$ coefficient products, each on
rationals of bit length $\poly(d_K,\beta_K)$, hence bit cost $\poly(d_K,\beta_K)$. Inversion uses the
extended Euclidean algorithm in $\mathbb{Q}[\theta]$, of the same order. The zero test is the exact
test on coordinates; no tolerance or numerical approximation enters.

\emph{(iii) The real embedding.} The element $\mu$ is a real root of $\chi_M$, isolated by exact real
root isolation (S2) into a rational interval $I$ that contains no other root. The embedding
$\sigma\colon K\to\mathbb{R}$ sends $\theta$ to that root; it is used only to decide \emph{which} real
root of $f_K$ a direction corresponds to, never to evaluate arithmetic, which stays in $K$. If $f_K$
has several real roots, they are Galois-conjugate and the corresponding instances are conjugate as well,
so the verdict is the same for all of them and the choice of $I$ does not affect it.


The purpose of this step is to bound the height of every object the procedure constructs; it is used
in every cost estimate of Section~\ref{sec5:sub:main}.

\begin{claim}[Heights of the data, of the roots and of the field]\label{sec5:lem:height}\label{sec5:lem:charpoly}\label{sec5:lem:roots}\label{sec5:lem:eigvec}
Let the entries of the input have bit length at most $\tau$.
\begin{enumerate}[label=(\alph*),leftmargin=*]
\item \emph{Compression.} A primitive integer basis $N$ of $K_0$ and the compression
$M=(N^{\mathsf T}N)^{-1}N^{\mathsf T}AN$ can be computed in bit complexity $\poly(n,\tau)$, and the
entries of both have bit length $\poly(n,\tau)$.
\item \emph{Characteristic polynomial.} After clearing denominators of $\chi_M(t)=\det(tI-M)$ and
taking the primitive part, the coefficients are integers of bit length $\poly(n,\tau)$, computable in
bit complexity $\poly(n,\tau)$.
\item \emph{Roots and field parameters.} Every root $\mu$ of $\chi_M$ satisfies
$[\mathbb{Q}(\mu):\mathbb{Q}]\le k\le n$ and $h(\mu)\le\poly(n,\tau)$, where $h$ is the logarithmic
absolute height; the real roots can be isolated into rational intervals of any prescribed accuracy in
bit complexity $\poly(n,\tau)$. The defining polynomial of Definition~\ref{sec5:def:model} satisfies
$d_K\le k\le n$ and $\beta_K=\poly(n,\tau)$, and is obtained in bit complexity $\poly(n,\tau)$.
\item \emph{Eigenvectors and fibre bases.} Under Definition~\ref{sec5:def:gsimple}, for every real
eigenvalue $\mu$ the space $E_\mu$ has a primitive vector whose coordinates in $K$ have size
$\poly(n,\tau)$, and bases of $E_\mu^0$, $U_v^\perp$, $D_v=K_0\cap U_v^\perp$ and
$\mathcal{Q}_v/\overline{V}_{0,v}$ can be computed over the same field $K$, in bit complexity
$\poly(n,\tau)$, with element sizes $\poly(n,\tau)$.
\end{enumerate}
\end{claim}

\begin{proof}
(a) Clear denominators, run fraction-free (Bareiss) elimination, normalize primitively and compress;
Hadamard's inequality controls every intermediate minor. (b) The coefficients are signed sums of
principal minors. (c) The degree is at most $\deg\chi_M=k$, and Cauchy's bound turns the coefficient
size of $\chi_M$ into a bound on $h(\mu)$; real root isolation is polynomial-time for integer
polynomials. The factor $f_K$ of the integer polynomial $\chi_M$ has height bounded by the height and
degree of $\chi_M$, which gives $\beta_K=\poly(n,\tau)$. (d) Each of the four spaces is a kernel, image,
sum or quotient of matrices of size at most $n$ over $K$, so fraction-free elimination over $K$
produces its basis, the coordinates being ratios of minors. The explicit bounds and their derivations
are in Appendix~\ref{appBC}, Claims~\ref{appBC:lem:norm}--\ref{appBC:lem:bases}.
\end{proof}

\begin{remark}[the bases must be computed over $K$]\label{sec5:rem:sameK}
These bases must be computed over $K=\mathbb{Q}(\mu)$, not over $\mathbb{Q}$: when $\mu\notin\mathbb{Q}$
the space $E_\mu$ is a line over $K$ but its trace over $\mathbb{Q}$ is a sum of conjugate lines, so a
computation over $\mathbb{Q}$ would split one direction into $[K:\mathbb{Q}]$ of them. Direction
enumeration and fibre construction therefore take place in the same field.
\end{remark}


The purpose of this step is to count charts, minors and parameters; it is used in
Theorem~\ref{sec5:thm:poly}.

\begin{claim}[Size of the charts]\label{sec5:sub:chart}\label{sec5:lem:chart}
Fix a real eigenvalue $\mu$, a direction $[v]\in\mathbb{P}(E_\mu^0)$ and an affine chart of the fibre
$\Gr(s-\delta,\mathcal{Q}_v/\overline{V}_{0,v})$, with $s=q-k+1$,
$\delta=\dim V_0-\dim(V_0\cap K_0)$, $V_0=\row L_0$, and $\widetilde r:=r(v)-\delta$ the rank of the
quotient bundle (unrelated to the free block $R$, which no longer appears). A chart is given by a
symbolic $d_{\mathrm{chart}}\times(\widetilde r-d_{\mathrm{chart}})$ matrix $X$. Then:
\begin{enumerate}[label=(\roman*),leftmargin=*]
\item $d_{\mathrm{chart}}=(s-\delta)(r(v)-s)\le p^{2}$;
\item every chart row is affine in $X$, and every maximal minor has degree $\le\omega$ in $X$ and in
$\lambda$;
\item the number of maximal minors is $\binom{q}{\omega}\binom{k}{\omega}=k\binom{q}{s}=n^{O(p)}$ for
fixed $p$;
\item each fibre is covered by $\binom{\widetilde r}{s-\delta}=O(1)^{p}$ charts;
\item there are at most $k\le n$ eigenvalues, each contributing at most one direction.
\end{enumerate}
\end{claim}

\begin{proof}
Each item is a count: $x\mapsto x(\widetilde r-x)$ is maximal at $x=\widetilde r/2$ with
$\widetilde r\le\rho\le2p$; the chart rows are affine in $X$ and the pencil is linear in $\lambda$;
$\binom{k}{\omega}=k$ because $\omega=k-1$, and $\binom{q}{s}\le\sum_{j\le2p}\binom{n}{j}$ because
$s\le\rho\le2p$. The counts and the dense-size bound $D=n^{O(p^{2})}$ are in Appendix~\ref{appBC},
Claim~\ref{appBC:lem:chart}.
\end{proof}


The purpose of this step is to bound the cost of the gcd the procedure computes; it is used in
Theorem~\ref{sec5:thm:poly}.

\begin{claim}[Cost of the gcd]\label{sec5:sub:subres}\label{sec5:lem:subres}\label{sec5:lem:multigcd}
Let $K$ be a number field in the representation of Definition~\ref{sec5:def:model}, with parameters
$d_K\le n$ and $\beta_K=\poly(n,\tau)$. Let $f,g\in K[X][\lambda]$ with $\deg_\lambda\le\omega$ and
$\deg_X\le\omega$ in each of the $d$ parameters, $d\le p^{2}$, and coefficient size at most $\beta$ (in
the sense of Definition~\ref{sec5:def:model}). Put
\begin{equation}\label{sec5:eq:sizes}
\begin{gathered}
\beta_S:=2\omega\bigl(\beta+\omega\log_2(2\omega)\bigr)+\poly(d_K,\beta_K),\\
D:=\binom{d+2\omega^{2}}{d}.
\end{gathered}
\end{equation}
Then:
\begin{enumerate}[label=(\alph*),leftmargin=*]
\item \emph{Sizes.} $\operatorname{Syl}(f,g)$ is a $(2\omega)\times(2\omega)$ matrix; the chain has at
most $\omega+1$ subresultants $S_0,\dots,S_\omega$; each $S_j$ lies in $K[X][\lambda]$, has
$\deg_\lambda S_j\le j$, $\deg_X S_j\le 2\omega^{2}$, and coefficient size at most $\beta_S$.
\item \emph{Gcd.} The primitive part of the gcd over $K(X)[\lambda]$ is $S_{j^{*}}$ up to a unit of
$K[X]$, where $j^{*}=\max\{j:S_j\neq0\}$; hence it has $\deg_\lambda=j^{*}$ and coefficient size at
most $\beta_S$. In particular the test ``$\deg_\lambda\gcd>0$'' is the test $j^{*}>0$.
\item \emph{Cost.} The chain is computed by the fraction-free subresultant algorithm in bit complexity
$\poly(D,\beta_S)$, and $D=n^{O(p^{2})}$ once $p$ is fixed.
\item \emph{Several polynomials.} For $m_1,\dots,m_\ell\in K[X][\lambda]$ with $\ell=n^{O(p)}$ and the
same degree and size bounds, the iterated gcd $\gcd(\dots\gcd(m_1,m_2),\dots,m_\ell)$ costs
$\poly(\ell,D,\beta_S)$ and has coefficient size at most $\beta_S$ at every step. The size is
\emph{bounded}, not necessarily decreasing: the primitive gcd is a factor of the current input, so it
satisfies the same factor-height bound as that input, and every intermediate subresultant satisfies the
subresultant bound of item (a). The number of non-trivial steps is at most $\omega$, since the
$\lambda$-degree strictly decreases whenever the gcd is non-trivial.
\end{enumerate}
\end{claim}

\begin{proof}
(a) Each entry of the Sylvester matrix is a coefficient of $f$ or $g$, of $X$-degree $\le\omega$ and
size $\beta$; a $(2\omega)\times(2\omega)$ determinant has $X$-degree at most $2\omega^{2}$ and, by
Hadamard's inequality applied to the $2\omega$ rows, coefficient size at most $\beta_S$ as in
\eqref{sec5:eq:sizes}. (b) The subresultant chain satisfies $\deg_\lambda\gcd=\max\{j:S_j\neq0\}$ and
$S_j\in K[X][\lambda]$; the primitive gcd divides the subresultant, hence obeys the same bound. (c) The
dense representation of a polynomial with $\lambda$-degree $\le\omega$ and $X$-degree $\le 2\omega^{2}$
in $d$ variables has $D=\binom{d+2\omega^{2}}{d}$ coefficients, and the algorithm performs $O(\omega)$
pseudo-remainder steps on such objects. (d) By Gauss's lemma over the UFD $K[X]$, the primitive gcd is a
factor of $m_1$ in $K[X][\lambda]$, so by the factor-height bound (S5) its coefficient size is at most
$\beta_S$; the iteration therefore stays within the bound rather than shrinking it. The counts are in
Appendix~\ref{appBC}, Claims~\ref{appBC:prop:iter} and~\ref{appBC:prop:cost}, and the failure of
the naive Euclidean algorithm is shown there on an explicit example.
\end{proof}


With the sizes bounded, the theorem is a bookkeeping of the steps.


\subsection{The fixed-\texorpdfstring{$p$}{p} theorem}\label{sec5:sub:main}

\begin{thm}[Fixed $p$, geometrically simple spectrum]\label{sec5:thm:poly}
Let $p$ be fixed. On inputs with geometrically simple real spectrum
(Definition~\ref{sec5:def:gsimple}), deciding feasibility in the hyperplane layer ($c=1$) can be
carried out in $\poly(n,\tau)$, concretely $n^{O(p^{2})}$, bit operations, using exact arithmetic
throughout.
\end{thm}

\begin{proof}
The steps and their costs are as follows.
\begin{enumerate}[label=(\arabic*),leftmargin=*]
\item \emph{Compression:} $\poly(n,\tau)$ by Claim~\ref{sec5:lem:height}(a).
\item \emph{Factorization, field and real roots:} factor $\chi_M$ into irreducibles by (S1), build $K$
as in Definition~\ref{sec5:def:model}, and isolate the real roots by (S2), all $\poly(n,\tau)$; by
Claim~\ref{sec5:lem:height}(c) each real root $\mu$ satisfies $[K:\mathbb{Q}]\le n$,
$h(\mu)=\poly(n,\tau)$ and $\beta_K=\poly(n,\tau)$.
\item \emph{Directions:} under geometric simplicity $E_\mu$ is one-dimensional and $E_\mu^0$ is a
subspace of it, so each $\mu$ yields at most one candidate direction; the eigenvector and the test for
membership in $E_\mu^0$ cost $\poly(n,\tau)$ by Claim~\ref{sec5:lem:height}(d).
\item \emph{Fibres:} bases of $D_v$, $U_v^\perp$ and $\mathcal{Q}_v/\overline{V}_{0,v}$, again by
Claim~\ref{sec5:lem:height}(d).
\item \emph{Charts:} by Claim~\ref{sec5:lem:chart} there are $\le2^{2p}$ charts per fibre,
$n^{O(p)}$ maximal minors, and degrees $\le\omega\le n$, while $d_{\mathrm{chart}}\le p^{2}$ is a
constant because $p$ is fixed.
\item \emph{Decision:} on each chart compute the subresultant gcd by
Claim~\ref{sec5:lem:subres}(a)--(d), at cost $\poly(D,\beta_S)$ with $D=n^{O(p^{2})}$ and
$\beta_S=\poly(n,\tau)$ by \eqref{sec5:eq:sizes}; $D$ is polynomial in $n$ since $p$ is fixed.
\item \emph{Decision rule:} by the full-rank leading-coefficient lemma of Section~\ref{sec4},
$\rank H\equiv\omega$ on every chart, so ``the gcd over $K(X)[\lambda]$ has degree zero'' is equivalent
to ``every point of the chart is good''. Hence the layer is feasible if some chart is generically
good, and infeasible if every chart is entirely bad.
\end{enumerate}
There are at most $k\cdot2^{2p}$ pieces and each term is either $\poly(n,\tau)$ or $n^{O(p^{2})}$, so
the total is $n^{O(p^{2})}$; the arithmetic in $K$ costs $\poly(d_K,\beta_K)$ per operation, which is
absorbed in the same bound.
\end{proof}

\paragraph{Promise class} The theorem is a \emph{promise-class} statement about \emph{cost}, not about
decidability: the procedure it prices is Algorithm~\ref{sec4:algo:decide}, which is defined on every input
and, as Theorem~\ref{sec4:prop:correct} states, decides each input on which it terminates; no step of
Section~\ref{sec4} assumes geometric simplicity, and the algorithm has no test that could report an input
as outside its scope. The hypothesis of Definition~\ref{sec5:def:gsimple} enters at one point only, in
step 3 of the proof of Theorem~\ref{sec5:thm:poly}, where it makes the direction parameter a finite list
instead of a projective eigenspace; that is what turns the estimate into $n^{O(p^{2})}$. On an input that
violates the promise the same algorithm still returns a correct verdict, but the dense size becomes
$n^{O(d_\mu+p^{2})}$ (the closing paragraph of this section), so no bound in $p$ alone follows and the
theorem asserts nothing about the running time outside the promise.

Five facts are quoted and not reproved: \emph{(S1)} factorization of integer polynomials into
irreducibles in bit complexity $\poly(N,\beta)$ \cite{vzgg03}; \emph{(S2)} real root isolation into
pairwise disjoint rational intervals at the same cost \cite{bpr06}; \emph{(S3)} arithmetic in a number
field $K=\mathbb{Q}[\mu]/(f)$ of degree $d$ and height $h$ in $\poly(d,h)$ bits with an exact zero
test \cite{lan02,vzgg03}; \emph{(S4)} the subresultant chain yields the primitive part of the gcd over
the fraction field, every intermediate quantity being a Sylvester minor with polynomially controlled
bit length \cite{vzgg03}; \emph{(S5)} the height of a factor of a polynomial is bounded by a polynomial
in the degree and the height of that polynomial (the factor-height, or Mignotte, bound), so the gcd of
a family of polynomials is no larger, in the sense of coefficient size, than the polynomials it divides
\cite{vzgg03,lan02}. Only the existence of polynomial-time algorithms is used, not the constants.


Every estimate the theorem needs is now in place, so the promise it carries can be examined.


\paragraph*{The repeated case}\label{sec5:sub:scope}

\textbf{Why the promise cannot be dropped \emph{for free}.} If some real eigenvalue has
$d_\mu=\dim E_\mu\ge2$, the base at $\mu$ has dimension $d_\mu-1$ and the number of chart parameters
becomes $(d_\mu-1)+(s-\delta)(r(v)-s)$; the second term is still $\le p^{2}$, but the first grows with
$n$, so the dense size becomes $n^{O(d_\mu+p^{2})}$, which is not polynomial for fixed $p$. Three
things must be kept apart: the geometry survives (the parameterization holds for arbitrary $d_\mu$),
the exact decision survives (the chartwise computation is correct for arbitrary $d_\mu$), and only the
\emph{complexity bound} fails.

\textbf{What the example of Appendix~\ref{app:example} proves.} It proves that inside a
repeated eigenspace the verdict can depend on the projective direction, so that selecting one
eigenvector per eigenvalue is not a valid reduction. It does \emph{not} prove that the repeated case is
computationally hard, and it does not prove that geometric simplicity is necessary for polynomial-time
solvability: no lower bound of any kind is asserted here. The complexity of the general
(repeated-spectrum) case is left open.

\textbf{Other cases not covered.} Layers with $c\ge2$ admit no corresponding direction--fibre
decomposition and are outside the scope of this paper. The design matrix is assumed to have full row
rank, which forces $\rank L_0=r_0$ as a necessary condition for feasibility and is used as a zero-cost
precondition. Real points cause no difficulty, every parameter chart being a rational real chart with
Zariski-dense real points.


What the instance shows, and what it does not, is stated in Appendix~\ref{app:example} as well.


\section{Discussion and conclusion}

This section states what is and is not claimed and records the evidence behind the claims; the
appendices carry the proofs of the estimates and the certified counterexample.

\subsection{Limitations and what is not claimed}\label{sec6:sub:conclusion}


Three statements of different strength are made, and they are kept apart throughout: (i) an exact
decision procedure for the whole layer, in exact arithmetic and with no complexity claim
(Section~\ref{sec4}); (ii) a polynomial bound on the \emph{cost} of that same procedure on the
geometrically simple subclass, a promise-class statement about running time rather than about which
inputs are decided (Section~\ref{sec5}); (iii) a counterexample refuting one specific reduction
(Appendix~\ref{app:example}). Nothing is claimed about the complexity of the repeated case.

\textbf{Limitations.} The conclusion does \emph{not} prove that the repeated-eigenspace case is
computationally hard (it refutes a particular reduction, not the existence of a general algorithm); it
does \emph{not} show that geometric simplicity of the real spectrum is necessary for the polynomial bound
of Theorem~\ref{sec5:thm:poly} (that hypothesis controls the dimension of the direction parameter, and the
example only shows that this dimension cannot be ignored); and it asserts nothing about the untested
directions of $E_1$.


\subsection{Reproducibility and evidence}

\begin{table*}[!t]
\centering
\caption{Where each statement of the paper is verified. \emph{Literal}: the Lean statement is the
statement of the paper. \emph{Abstract core}: the Lean statement is the linear-map or submodule
version, and the translation to matrices is paper-level. \emph{Ingredient}: one step of the proof,
not the statement. Statements not listed are paper-level throughout.}
\label{tab:evidence}
\footnotesize
\begin{tabular}{@{}p{0.43\textwidth}>{\raggedright\arraybackslash}p{0.15\textwidth}>{\raggedright\arraybackslash}p{0.36\textwidth}@{}}
\toprule
Statement & Lean evidence & Formal name \\
\midrule
Direction theorem (Theorem~\ref{sec3:thm:direction}), $\varphi_{Av}=\mu\varphi_v\Rightarrow Mv=\mu v$
& abstract core & \texttt{compression\_eq\_of\_proportional} \\
Direction theorem (Theorem~\ref{sec3:thm:direction}), $Mv=\mu v\Rightarrow\varphi_{Av}=\mu\varphi_v$,
and the equivalence of the two & abstract core & \texttt{proportional\_of\_compression},
\texttt{direction\_iff} \\
Structure of the layer (Claim~\ref{sec3:lem:direction}(ii)): membership in $U_v^\perp$ &
ingredient & \texttt{mem\_orthogonal\_span\_pair\_iff} \\
Prescribed rows (Claim~\ref{sec3:lem:prescribed}(i)): $V_0\subseteq U_v^\perp$ &
abstract core & \texttt{le\_orthogonal\_span\_pair} \\
Prescribed rows (Claim~\ref{sec3:lem:prescribed}(i)): $V_0\cap D_v=V_0\cap K_0$, so $\delta$ is
independent of $v$ & abstract core & \texttt{inf\_inf\_orthogonal\_eq} \\
Fibre: $\dim D_v=k-1$ (Claim~\ref{sec3:lem:direction}(iii),
Claim~\ref{sec3:lem:constantrank}) & abstract core & \texttt{Dv\_eq\_ker},
\texttt{finrank\_Dv} \\
Avoidance of finitely many proper hyperplanes (the linear model of the generic-point step of
Claim~\ref{sec4:thm:generic}) & ingredient & \texttt{exists\_apply\_ne\_zero},
\texttt{exists\_forall\_apply\_ne\_zero} \\
\bottomrule
\end{tabular}
\end{table*}

\textbf{Reproducibility.} The matrices $A$, $C$, $N$, $M$, $A_{21}$ and the characteristic polynomial
are given in full in Appendix~\ref{sec6:sub:instance}; the six decisions, the divisibility certificates
and the chart-point certificate are exact rational computations whose expanded forms, together with a
short reproduction recipe, are given in Appendix~\ref{app:data}.

\emph{What these certificates establish.} Identities \eqref{appA:eq:divcert} and
\eqref{appA:eq:bezout} are finite rational identities: expanding either side and cancelling verifies
them in any exact arithmetic system, with no algebraic extension and no rounding. Together with the
full-rank leading-coefficient lemma of Section~\ref{sec:decision} they give the two conclusions used in
Appendix~\ref{app:example}: the fibres over the five non-eigen directions are entirely bad, and
the fibre over $(1\!:\!1)$ contains a feasible design.

\textbf{Lean evidence.} The statements of Section~\ref{sec3} that concern the layer itself --- the
direction theorem, the prescribed rows, and the dimension of the fibre --- are also formalized in
Lean~4 with Mathlib, as are the inner-product-space facts on which they rest;
Table~\ref{tab:evidence} lists them statement by statement, with the grade of the formalization and
the declaration names, in the format used by the companion papers. The grades separate what is
formalized as stated (\emph{literal}) from what is formalized at the level of linear maps, submodules
and inner product spaces (\emph{abstract core}), the translation to the matrices $M$, $A_{21}$, $L_0$
being paper-level, from what is one step of a proof (\emph{ingredient}). Every declaration is checked
by the kernel against only (\texttt{propext}, \texttt{Classical.choice}, \texttt{Quot.sound}). The
statements that have no formalization --- the closedness of the bad locus, the function-field gcd
criterion of Claim~\ref{sec4:thm:generic}, the bundle and Grassmannian description of the fibre in
Claim~\ref{sec3:lem:constantrank}, and the algorithm with its correctness and complexity bounds --- are
paper-level throughout.

Everything above is exact rational arithmetic whose expanded forms are in Appendix~\ref{app:data}, and
the proofs behind every estimate of Section~\ref{sec5} are in Appendix~\ref{appBC}; nothing in this
paper rests on a numerical computation.


\appendices
\small
\section{Proofs of the estimates}\label{appBC}\label{appB}\label{app:subresultant}


This appendix proves the bounds that Section~\ref{sec5} states in $\poly$ form, recording the explicit
estimates. Notation is that of Sections~\ref{sec:prelim} and~\ref{sec5}; $\tau$ is the encoding length
of $(A,C,L_0,q)$; $\beta(x)$ is the largest bit length of the numerator and denominator of a rational
$x$, extended entrywise to matrices; logarithms are to base $2$; constants are not optimized.

Every field that occurs below is a number field in the representation of
Definition~\ref{sec5:def:model}: $K=\mathbb{Q}[\theta]/(f_K)$ with $f_K\in\mathbb{Z}[\theta]$
primitive of degree $d_K\le n$ and height $\beta_K=\poly(n,\tau)$, elements represented by coordinate
vectors in the power basis, and the \emph{size} of an element meaning the largest bit length of those
coordinates. Thus a statement such as ``height at most $\beta$'' always refers to that representation,
and every bound below is uniform in the field because $d_K$ and $\beta_K$ are themselves bounded by
$\poly(n,\tau)$ (Claim~\ref{appBC:lem:bases}).


\subsection{Proof of Claim~\ref{sec5:lem:height}: heights}\label{appBC:sub:norm}


Clear denominators in $[C;CA]$, multiplying by the least common multiple $D_0$ of the denominators:
this gives an integer matrix $\widehat W$ with $\beta(\widehat W)\le\tau+2pn\tau=:\beta_W$, because
$\beta(D_0)\le2pn\tau$. Its kernel is computed by Bareiss elimination, whose intermediate quantities
are minors, and Hadamard's inequality bounds a $j\times j$ minor by
$\log_2|\det|\le\frac{j}{2}\log_2 j+j\beta_W\le\frac{\rho}{2}\log_2\rho+\rho\beta_W$ for
$j\le\rho$. Clearing denominators in the resulting coordinates and dividing each column by the gcd of
its entries (the \emph{primitive normalization}) adds no bits and keeps $\det(N^{\mathsf T}N)$ an
integer, without which a common rational factor of $N$ would inflate both parts of every entry of $M$
\cite{hj13,vzgg03}.


The purpose of this step is the explicit size of the compressed data; it is used in
Claim~\ref{appBC:lem:charpoly} and in Section~\ref{sec5}.

\begin{claim}[explicit form of Claim~\ref{sec5:lem:height}(a)]\label{appBC:lem:norm}\label{appBC:lem:comp}\label{appBC:cor:comp}
Let $\beta_W=\tau+2pn\tau$. Then
\begin{enumerate}[label=(\roman*),leftmargin=*]
\item $\beta(N)\le2\bigl(\frac{\rho}{2}\log_2\rho+\rho\beta_W\bigr)$, and $N$ is computable in bit
complexity $\poly(n,\tau)$;
\item $\beta(N^{\mathsf T}N)\le2\beta(N)+\log_2 n$ and
$\beta(N^{\mathsf T}AN)\le2\beta(N)+\tau+2\log_2 n$;
\item the entries of $M$ satisfy
$\beta(M)\le(k-1)\bigl(\tfrac12\log_2 k+\beta(N^{\mathsf T}N)\bigr)+\beta(N^{\mathsf T}AN)+\beta\bigl(\det(N^{\mathsf T}N)\bigr)$,
hence $\beta(M)\le c_1(n^{3}\tau+n^{2}\log_2 n)$ for a constant $c_1$.
\end{enumerate}
\end{claim}

\begin{proof}
(i) Each step has a polynomial arithmetic count in $n$ and acts on integers of $\poly(n,\tau)$ bits.
(ii) The entries of $N^{\mathsf T}N$ (respectively $N^{\mathsf T}AN$) are sums of $n$ (respectively
$n^{2}$) terms of size at most $2^{2\beta(N)}$ (respectively $2^{2\beta(N)+\tau}$). (iii) $N$ has full
column rank, so $N^{\mathsf T}N$ is integral and positive definite and $\det(N^{\mathsf T}N)\ge1$; the
numerator of $(N^{\mathsf T}N)^{-1}=\operatorname{adj}(N^{\mathsf T}N)/\det(N^{\mathsf T}N)$ is thus
controlled by the adjugate bound alone, whose entries are $(k-1)$-minors bounded by Hadamard
\cite{hj13,vzgg03}.
\end{proof}


The purpose of this step is the explicit size of the field data and of the bases; it is used in
Claim~\ref{sec5:lem:height} and in Claim~\ref{sec5:lem:chart}.

\begin{claim}[explicit form of Claim~\ref{sec5:lem:height}(b)--(d)]\label{appBC:lem:charpoly}\label{appBC:lem:root}\label{appBC:cor:root}\label{appBC:lem:bases}\label{appBC:sub:roots}
Let $K=\mathbb{Q}(\mu)$ be represented as in Definition~\ref{sec5:def:model}, with
$d_K=[K:\mathbb{Q}]\le n$ and $\beta_K=\beta(f_K)$.
\begin{enumerate}[label=(\roman*),leftmargin=*]
\item after clearing denominators and taking the primitive part,
$\beta(\chi_M)\le k\bigl(\beta(M)+\tfrac12\log_2 k\bigr)+\beta(M)+k=\poly(n,\tau)$;
\item every root $\mu$ of $\chi_M$ satisfies $d_K\le k\le n$ and
$h(\mu)\le\beta(\chi_M)+\tfrac12\log_2(k+1)+1$, so $\beta_K\le\beta(\chi_M)+\tfrac12\log_2(k+1)+1$;
the real roots are separable into pairwise disjoint rational intervals in bit complexity
$\poly(n,\tau)$ \cite{bpr06};
\item $E_\mu=\ker(M-\mu I)$ has over $K$ a primitive vector $a$ with
$\beta(a)\le2\bigl(\tfrac k2\log_2 k+k(\beta(M)+\beta_K)\bigr)$, bases of $E_\mu^0$,
$U_v^\perp$, $D_v=K_0\cap U_v^\perp$ and $\mathcal{Q}_v/\overline{V}_{0,v}$ are computable over $K$
with element sizes $\poly(n,\tau)$, and the bit complexity of all these computations is
$\poly(n,\tau)$.
\end{enumerate}
\end{claim}

\begin{proof}
(i) $\chi_M=\sum_i(-1)^{i}\sigma_i(M)t^{k-i}$ with $\sigma_i$ a sum of principal $i$-minors, each
bounded by $(\sqrt i\,\|M\|_\infty)^{i}$; summing $\binom ki\le2^{k}$ terms and clearing denominators
adds $\poly(n,\tau)$ bits, and Faddeev--LeVerrier or interpolation gives a polynomial arithmetic count.
(ii) Cauchy's bound gives $|\mu|\le1+\max_i|a_i/a_k|\le2^{\beta(\chi_M)+1}$, so each conjugate
satisfies $\log^{+}|\sigma(\mu)|\le\beta(\chi_M)+1$, and dividing by the degree gives the height bound
\cite{hj13}. (iii) The vector $a$ spans $\ker(M-\mu I)$ over $K$ and its coordinates are ratios of
$(k-1)$-minors over $K$ to a nonzero principal one, bounded by Hadamard together with
$\beta(M),\beta_K$; each of the four spaces is a kernel, image, sum or quotient of matrices over $K$
of size at most $n$, and all inner products lie in $K$; the complexity follows from $d_K\le n$, a
polynomial arithmetic count, and arithmetic in $K$ at cost $\poly(d_K,\beta_K)$ bits
\cite{lan02,vzgg03}.
\end{proof}

\begin{remark}[why the computation cannot be carried out over $\mathbb{Q}$]\label{appBC:rem:mustK}
When $\mu\notin\mathbb{Q}$ the space $E_\mu$ is a line over $K$ but not over $\mathbb{Q}$: its trace
over $\mathbb{Q}$ is a direct sum of $d$ conjugate lines. Computing the kernel over $\mathbb{Q}$ would
split one direction into $d$ conjugate directions and change both the fibre and the pencil, so
direction enumeration and fibre construction must take place in $K$. Since $\mu$ is never turned into a
number, nothing here grows with the distance from $\mu$ to $\mathbb{Q}$.
\end{remark}


\subsection{Proof of Claim~\ref{sec5:lem:chart}: sizes of the charts}\label{appBC:sub:chart}


Fix a real eigenvalue $\mu$ and a direction $[v]\in\mathbb{P}(E_\mu^0)$; the fibre is
$\Gr(m,\mathcal{Q}_v/\overline{V}_{0,v})$ with $m=s-\delta$, and $\widetilde r:=r(v)-\delta$ is the
rank of the quotient bundle, where $r(v)=\rho$ for eigen directions and $\rho-1$ otherwise.


The purpose of this step is the explicit chart count and dense size; it is used in
Theorem~\ref{sec5:thm:poly}.

\begin{claim}[explicit form of Claim~\ref{sec5:lem:chart}]\label{appBC:lem:chart}\label{appBC:rem:poly}
With $d_{\mathrm{chart}}=m(\widetilde r-m)$:
\begin{enumerate}[label=(\roman*),leftmargin=*]
\item $d_{\mathrm{chart}}\le\widetilde r^{2}/4\le p^{2}$;
\item an $\omega\times\omega$ maximal minor has degree at most $\omega$ in $X$ and in $\lambda$;
\item the number of minors is
$\binom q\omega\binom k\omega=k\binom qs\le n\sum_{j\le2p+1}\binom nj=n^{O(p)}$;
\item the chart cover of $\Gr(m,\widetilde r)$ has $\binom{\widetilde r}{m}\le2^{\widetilde r}\le2^{2p}$
elements, with at most $k\le n$ pieces, hence at most $n\cdot2^{2p}$ charts;
\item the dense size of a maximal minor, and of every subresultant below, is
\begin{equation}\label{appBC:eq:denseminor}
D=\binom{d_{\mathrm{chart}}+2\omega^{2}}{d_{\mathrm{chart}}},\qquad d_{\mathrm{chart}}\le p^{2},
\end{equation}
so $D\le\binom{p^{2}+2n^{2}}{p^{2}}=n^{O(p^{2})}$, a polynomial in $n$ once $p$ is fixed, the exponent
involving $p^{2}$ only.
\end{enumerate}
\end{claim}

\begin{proof}
$\widetilde r\le\rho\le2p$ and $x\mapsto x(\widetilde r-x)$ is maximal at $x=\widetilde r/2$, giving
(i); a chart row is affine in $X$ and the pencil is linear in $\lambda$, giving (ii);
$\binom k\omega=k$ because $\omega=k-1$, and $\binom qs\le\sum_{j\le2p}\binom nj$ because
$s\le\rho\le2p$, giving (iii); (iv) uses $\sum_j\binom{\widetilde r}j=2^{\widetilde r}$; (v) uses
$\binom{d+\omega}d\le\bigl(e(d+\omega)/d\bigr)^{d}$ with $d\le p^{2}$ constant and $\omega\le n$.
\end{proof}

A resultant of two $\lambda$-polynomials of degree $\omega$ is a Sylvester determinant of order
$2\omega$ with $X$-degree at most $\omega$ and height at most $\beta$, hence of bit length
$\poly(\omega,\beta)$; the criterion built on it is only \emph{sufficient}, failure being no certificate
that a chart is bad.


\subsection{Proof of Claim~\ref{sec5:lem:subres}: the subresultant chain}\label{appBC:sub:subres}


\paragraph{Why naive Euclid fails}\label{appBC:sub:naive}
Division with remainder directly in $K(X)[\lambda]$ is the wrong procedure: from the second step on the
objects divided are rational functions, and each step multiplies the reciprocal of the leading
coefficient of the divisor into the next denominator, so denominator degrees can double at every step.
Take $K=\mathbb{Q}$ and one parameter $X$, with
\[
f=\lambda^{3}+X,\qquad g=\lambda^{2}+X\lambda+1 ,
\]
coprime over $\mathbb{Q}(X)[\lambda]$, and run the algorithm nevertheless. Dividing $f$ by $g$ gives
\[
q_1=\lambda-X,\qquad r_1=f-q_1g=(X^{2}-1)\lambda+2X ,
\]
with \emph{no} denominator; dividing $g$ by $r_1$ gives
\[
\begin{gathered}
q_2=\frac{\lambda}{X^{2}-1}+\frac{X^{3}-3X}{(X^{2}-1)^{2}},\\
r_2=g-q_2r_1=\frac{-X^{4}+4X^{2}+1}{(X^{2}-1)^{2}} ,
\end{gathered}
\]
a nonzero constant in $\lambda$, so $\gcd(f,g)=1$. The denominator jumps from none to degree $4$ in a
single step, and the numerator is exactly $\operatorname{Res}_\lambda(f,g)$; in general the quotient at
step $i$ carries $1/\operatorname{lc}_\lambda(r_{i-1})$, so denominators accumulate multiplicatively and
after $\omega$ steps can reach $2^{\Omega(\omega)}$. Direct Euclidean division over $K(X)[\lambda]$
therefore cannot bound coefficient size by $\poly(n,\tau)$.

\paragraph{The subresultant chain}
The correct tool is the \emph{subresultant}, a determinant computed inside the polynomial ring, so no
denominators arise. For $f=\sum_{i\le m}f_i\lambda^{m-i}$ and $g=\sum_{j\le n}g_j\lambda^{n-j}$ in
$K[X][\lambda]$, the Sylvester matrix $\operatorname{Syl}(f,g)$ is the $(m+n)\times(m+n)$ matrix of
shifted coefficient vectors with $\det=\operatorname{Res}(f,g)$, and for $0\le j<\min(m,n)$ the $j$-th
subresultant $S_j$ is the determinant obtained by deleting the last $j$ rows of each block and the last
$2j$ columns, so $S_0=\operatorname{Res}(f,g)$. Three properties matter: each $S_j$ lies in
$K[X][\lambda]$ (\emph{fraction-free}, the decisive advantage over Euclid), $\deg_\lambda S_j\le j$
(\emph{degree}), and with $j^{*}=\max\{j:S_j\neq0\}$ one has
$\deg_\lambda\gcd_{K(X)[\lambda]}(f,g)=j^{*}$ (\emph{gcd degree}), so $f,g$ are coprime if and only if
$S_0\neq0$. The sizes are as follows. The inputs satisfy
\begin{equation}\label{appBC:eq:sizesin}
\begin{gathered}
\deg_\lambda f,\ \deg_\lambda g\le\omega,\qquad
\deg_X f,\ \deg_X g\le\omega,\\
\text{coefficient size}\le\beta ,
\end{gathered}
\end{equation}
where $X$ carries $d\le p^{2}$ parameters and $\beta$ is measured in the representation of $K$. Hence
\begin{enumerate}[label=(\alph*),leftmargin=*]
\item $\operatorname{Syl}(f,g)$ has order $m+n\le2\omega$, and each of its entries is a coefficient of
$f$ or $g$, so of $X$-degree $\le\omega$ and size $\le\beta$;
\item the chain has $m+1\le\omega+1$ members $S_0,\dots,S_m$, and it is computed only up to the first
vanishing member, so its length is at most $\omega$;
\item each $S_j$ is a $(2\omega-j)\times(2\omega-j)$ determinant of such entries, hence
\begin{equation}\label{appBC:eq:sizesout}
\begin{gathered}
\deg_X S_j\le2\omega^{2},\\
\text{coefficient size}\le\beta_S,\\
\beta_S:=2\omega\bigl(\beta+\omega\log_2(2\omega)\bigr)+\poly(d_K,\beta_K),
\end{gathered}
\end{equation}
by Hadamard's inequality applied to the $2\omega$ rows;
\item the dense size of such an object is
\begin{equation}\label{appBC:eq:dense}
\begin{gathered}
D=\binom{d+2\omega^{2}}{d}=\binom{p^{2}+2\omega^{2}}{p^{2}},\\
D=n^{O(p^{2})}\ \text{for fixed } p,
\end{gathered}
\end{equation}
since $d\le p^{2}$ and $\omega\le n$; note that $D$ is polynomial in $n$ once $p$ is fixed, the
exponent depending on $p^{2}$ only.
\end{enumerate}
The cost of the fraction-free chain is $\poly(D,\beta_S)$, i.e.\ polynomial in the number of
coefficients and in their bit length, by (S4).


The purpose of this step is the cost of the iterated gcd; it is used in Claim~\ref{sec5:lem:subres}.

\begin{claim}[iteration and cost]\label{appBC:prop:iter}\label{appBC:prop:cost}\label{appBC:rem:pairwise}
Let $K$ be a number field in the representation of Definition~\ref{sec5:def:model}, and let
$m_1,\dots,m_\ell\in K[X][\lambda]$ satisfy \eqref{appBC:eq:sizesin} with $\ell=n^{O(p)}$. Then:
\begin{enumerate}[label=(\roman*),leftmargin=*]
\item $\gcd(m_1,\dots,m_\ell)=\gcd\bigl(\gcd(m_1,\dots,m_{\ell-1}),m_\ell\bigr)$, and iterating in the
order of the indices never increases the $\lambda$-degree or the $X$-degree;
\item the coefficient size of every iterated primitive gcd, and of every intermediate subresultant, is
at most $\beta_S$ of \eqref{appBC:eq:sizesout}, \emph{uniformly} over the iteration; it is bounded by
the input data, not decreased by the iteration;
\item the iterated algorithm computes the gcd within $\poly(\ell,D,\beta_S)$, and the test ``degree
positive'' agrees with ``all $\ell$ polynomials have a nontrivial common factor'';
\item a test based on \emph{pairwise} resultants would not do: for $f_1=\lambda(\lambda-1)$,
$f_2=\lambda(\lambda-2)$, $f_3=(\lambda-1)(\lambda-2)$ every pairwise resultant vanishes, yet
$\gcd(f_1,f_2,f_3)=1$.
\end{enumerate}
\end{claim}

\begin{proof}
(i) Associativity of the gcd; a gcd in $K[X][\lambda]$, taken primitive, divides its arguments in
$K[X][\lambda]$, so its $\lambda$-degree and its $X$-degree are bounded by those of the arguments and
never increase along the iteration. (ii) The primitive gcd of a step divides one of its arguments in
the UFD $K[X][\lambda]$ (Gauss's lemma), so by the factor-height bound (S5) its coefficient size is at
most the size bound attached to that argument, namely $\beta_S$ of \eqref{appBC:eq:sizesout}; the same
bound applies to every intermediate subresultant by item (a) of the description of the chain. The bound
is therefore uniform, but it is \emph{not} claimed to decrease from step to step: a gcd can be as
large, in coefficient size, as the input it divides. (iii) There are $\ell-1$ steps, each costing
$\poly(D,\beta_S)$ by (S4), and the total follows. (iv) Direct verification: $\lambda$ divides $f_1$,
$f_2$; $\lambda-1$ divides $f_1,f_3$; $\lambda-2$ divides $f_2,f_3$; no factor divides all three.
\end{proof}

\paragraph{Relation to the implementation}\label{appBC:sub:impl}
The complexity count of Section~\ref{sec5} is made for the subresultant chain, including its
multi-polynomial version; the implementation that produced the certificates of
Appendix~\ref{app:data} divides in $\mathbb{Q}(X)[\lambda]$ instead, which is feasible for the small
instances ($\omega\le3$, $d\le4$) with normalized denominators but costs what the example above
describes on large ones --- a limitation of the implementation, not of the theorem. Both compute the
\emph{same} object, the gcd in $K(X)[\lambda]$, whose degree-zero test is unaffected by removing the
content \cite{vzgg03}.

One further step passes from the gcd to a decision. Over a general family a positive gcd degree does
\emph{not} imply that every specialization has a common factor, since the factor may degenerate to a
constant ($X\lambda+1$ degenerates to $1$ at $X=0$); here the full-rank leading-coefficient lemma of
Section~\ref{sec4} supplies what is missing, because on a chart $\rank H\equiv\omega$, so at every
point some maximal minor has $\lambda$-degree exactly $\omega$ and the primitive common factor cannot
degenerate anywhere:
\[
\begin{gathered}
\bigl(\deg_\lambda\gcd>0\bigr)\ \wedge\ \bigl(\rank H\equiv\omega\bigr)\\
\Longrightarrow\quad \text{every point of the chart is bad}.
\end{gathered}
\]
\emph{Computing} the gcd (this appendix) and \emph{using} it to judge a whole chart
(Section~\ref{sec4}) are thus two separate matters.

\paragraph{Why the bad locus is closed}\label{appBC:sub:closed}
This is statement (a) of the Claim of Section~\ref{sec:prelim}, quoted there and not used elsewhere in the
paper; it is proved here for completeness. Fix a chart $\mathcal X$ and let $m_I$, with $I$ running over
the index pairs of Section~\ref{sec:prelim}, be its maximal minors, of $\lambda$-degree at most $\omega$ and
polynomial in $X$. Homogenize them to degree $\omega$ in
$(\lambda:\mu)$ and put
\[
Z=\bigl\{(X,(\lambda:\mu))\in\mathcal X\times\mathbb P^1:\ \tilde m_I(X,\lambda,\mu)=0\ \text{for all }I\bigr\}.
\]
$Z$ is closed in $\mathcal X\times\mathbb P^1$, and the second projection $\mathcal X\times\mathbb P^1\to
\mathcal X$ is proper, so its image is closed; equivalently, the image is the variety of the elimination
ideal of $Z$ in $K[X]$, which can be computed by resultants. On a chart $\rank H\equiv\omega$, so some
maximal minor of $H$ is nonzero and the corresponding $m_I$ has $\lambda$-degree exactly $\omega$; its
leading coefficient does not vanish, so $(\lambda:\mu)=(1:0)$ is not a point of $Z$ at any $X$ and the
image is exactly the bad locus $\mathcal B=\{X:\deg_\lambda\gcd_I m_I(X,\cdot)>0\}$. Hence $\mathcal B$ is
Zariski closed. Nothing here uses a companion paper, and the only hypothesis beyond polynomiality is the
non-vanishing of one leading coefficient, which is what $\rank H\equiv\omega$ provides.


\section{A repeated-eigenspace counterexample}\label{app:example}\label{sec:counterexample}\label{appA}\label{app:data}

This appendix presents a complete rational instance that refutes a natural simplification in the
hyperplane layer: \emph{taking, for each distinct real eigenvalue of the compression, one eigenvector as
the representative direction}, and deciding feasibility from that representative alone. In the instance
the compression has a two-dimensional real eigenspace, and different projective directions inside it
lead to \emph{different} answers: one direction is feasible, the others tested are infeasible. What is
refuted is exactly that reduction; the example does not show that the repeated case is computationally
hard, nor that geometric simplicity of the real spectrum is necessary for polynomial-time decidability
(Section~\ref{sec6:sub:conclusion}). All computations are exact rational arithmetic: infeasibility of a
direction is certified by a divisibility identity over the function field of the fibre, using the
maximal-minor test of \cite{companion1}, and feasibility of the good direction is certified at a single
rational chart point.

\subsection{Notation}\label{sec6:sub:notation}


We use the notation of the preceding sections. The appendix lives in the hyperplane layer $c=1$ and in
the transverse case $\dim(V\cap K_0)=k-c=k-1$, so $\ker H$ is one-dimensional, say
$\ker H=\mathbb{R}v$ with $v\in K_0$, every row of $\mathcal{L}$ annihilates $v$ and $Av$, and with
$U_v=\operatorname{span}\{v,Av\}$ the admissible row spaces obey $V\subseteq U_v^{\perp}$. The fibre
computation below uses only $V\subseteq U_v^{\perp}$ and $\dim(V\cap K_0)=k-1$: the inclusion gives
$V\cap K_0\subseteq D_v$, and the dimension gives $V\cap K_0=D_v$, whence $\dim V/D_v=q-k+1=s$. Write
\begin{equation}\label{sec6:eq:defs}
D_v=K_0\cap U_v^{\perp},\qquad \mathcal{Q}_v=U_v^{\perp}/D_v,\qquad s=q-k+1 .
\end{equation}
Since $L_0$ is empty, the admissible directions are exactly the real eigenvector directions of the
compression, and for fixed $[v]$ the fibre is $\mathbb{P}(\mathcal{Q}_v)$, of dimension $r(v)-1$ with
$r(v)=\dim\mathcal{Q}_v=\rho$ for an eigen direction ($Av\in\operatorname{span}v$) and $\rho-1$ for a
non-eigen direction. In this appendix $s=1$, so every fibre is a projective space and a single
computation over the function field of the fibre covers it entirely: no chart decomposition is needed.


\subsection{The instance}\label{sec6:sub:instance}


Take
\begin{equation}\label{sec6:eq:data}
n=8,\quad p=2,\quad q=4,\qquad
C=\begin{bmatrix}e_1\\ e_2\end{bmatrix},\qquad L_0\ \text{empty},
\end{equation}
and let the state matrix be
\begin{equation}\label{sec6:eq:A}
A=\begin{bmatrix}
0&0&1&0&0&0&0&0\\
0&0&0&1&0&0&0&0\\
1&-2&1&0&-1&1&1&-1\\
-1&2&1&0&2&-2&-1&2\\
0&0&0&0&1&0&0&0\\
0&0&0&0&0&1&0&0\\
0&0&0&0&0&0&2&0\\
0&0&0&0&0&0&0&3
\end{bmatrix}.
\end{equation}
Its lower-right block is $\operatorname{diag}(1,1,2,3)$, its upper-right block is
$\begin{bsmallmatrix}0&0&0&0\\0&0&0&0\\-1&1&1&-1\\2&-2&-1&2\end{bsmallmatrix}$, its lower-left block
vanishes, and its upper-left block has first two rows giving $CA=[e_3;e_4]$.

It follows that
\begin{equation}\label{sec6:eq:K0}
K_0=\operatorname{span}\{e_5,e_6,e_7,e_8\},
\qquad
N=\begin{bmatrix}
0&0&0&0\\0&0&0&0\\0&0&0&0\\0&0&0&0\\
1&0&0&0\\0&1&0&0\\0&0&1&0\\0&0&0&1
\end{bmatrix},
\end{equation}
hence
\begin{equation}\label{sec6:eq:invariants}
\rho=k=4,\qquad \omega=k-1=3,\qquad s=q-k+1=1 ,
\end{equation}
and, in the coordinates of $K_0$,
\begin{equation}\label{sec6:eq:M}
\begin{aligned}
M&=\pi_{K_0}A|_{K_0}=\operatorname{diag}(1,1,2,3),\\
\chi_M(x)&=(x-3)(x-2)(x-1)^2 .
\end{aligned}
\end{equation}
\begin{equation}\label{sec6:eq:A21}
A_{21}=\pi_{K_0^{\perp}}A|_{K_0}=
\begin{bmatrix}
0&0&0&0\\0&0&0&0\\-1&1&1&-1\\2&-2&-1&2\\
0&0&0&0\\0&0&0&0\\0&0&0&0\\0&0&0&0
\end{bmatrix}.
\end{equation}

The following identities hold in exact arithmetic: $[C;CA]=[e_1;e_2;e_3;e_4]$; $N^{\mathsf{T}}N=I_4$ and
$N^{\mathsf{T}}AN=\operatorname{diag}(1,1,2,3)$; $AN=NM+A_{21}$; and $\chi_M(x)=(x-3)(x-2)(x-1)^2$. The
instance is obtained from a larger example by keeping $C$ and the action of $A$ outside $K_0$ and
replacing only the action on $K_0$, which creates a two-dimensional real eigenspace while leaving
$K_0$, $\rho$, $k$ and $\omega$ unchanged.


\subsection{Direction geometry}\label{sec6:sub:direction}


The factor $(x-1)^2$ produces a two-dimensional eigenspace
$E_1=\ker(M-I)=\operatorname{span}\{e_1,e_2\}\subseteq\mathbb{R}^{4}$, in the coordinates of $K_0$
(corresponding to $e_5,e_6$ in the state space); the eigenspaces of the eigenvalues $2$ and $3$ are
one-dimensional. Directions in $E_1$ are written
\begin{equation}\label{sec6:eq:v}
v(c_1,c_2)=c_1e_5+c_2e_6,\qquad (c_1,c_2)\neq(0,0),
\end{equation}
so the direction space is the projective line $\mathbb{P}(E_1)\cong\mathbb{P}^{1}$.

\textbf{The eigen/non-eigen stratification.} For $c=(c_1,c_2,0,0)^{\mathsf{T}}\in E_1$,
\begin{equation}\label{sec6:eq:A21c}
\begin{aligned}
A_{21}c&=(c_2-c_1)\begin{bmatrix}1&-2&0&0&0&0&0&0\end{bmatrix}^{\mathsf{T}},\\
A_{21}c&=0\quad\iff\quad c_1=c_2 .
\end{aligned}
\end{equation}
Thus $E_1$ contains exactly \emph{one} eigen direction, $[1\!:\!1]$, and every other direction of $E_1$
is a non-eigen direction; this is consistent with the general description, since the eigen directions
inside $E_\mu$ form the linear subspace $\ker A_{21}$. For $[1\!:\!1]$ one checks directly that $Av=v$,
so $\dim U_v=1$.

\textbf{Fibres.} With $s=1$ and $\delta=0$ the fibre over a fixed direction is $\mathbb{P}^{2}$ for a
non-eigen direction ($r(v)=\rho-1=3$) and $\mathbb{P}^{3}$ for the eigen direction
($r(v)=\rho=4$): the eigen direction has $\dim U_v=1$ instead of $2$, hence $\dim U_v^{\perp}=7$
instead of $6$, while $\dim D_v=k-1=3$ is unchanged, so its fibre gains exactly one homogeneous
parameter. The two situations are compared in Table~\ref{sec6:tab:dims}.

\begin{table}[t]
\centering
\caption{Dimensions attached to a direction, for $n=8$ and $k=4$. The two columns differ because
$\dim U_v$ differs: an eigen direction places $Av$ inside $\operatorname{span}v$, so $U_v$ loses one
dimension and its annihilator gains one.}
\label{sec6:tab:dims}
\begin{tabular}{@{}lcc@{}}
\toprule
Quantity & Non-eigen & Eigen\\
\midrule
$\dim U_v$ & $2$ & $1$\\
$\dim U_v^{\perp}=n-\dim U_v$ & $6$ & $7$\\
$r(v)=\dim U_v^{\perp}-(k-1)$ & $3$ & $4$\\
Fibre $\mathbb{P}^{r(v)-1}$ & $\mathbb{P}^{2}$ & $\mathbb{P}^{3}$\\
\bottomrule
\end{tabular}
\end{table}

\textbf{Example: the direction $(1\!:\!0)$.} Here $v=e_5$ and
$Av=Ae_5=(0,0,-1,2,1,0,0,0)^{\mathsf{T}}$, so in the coordinates of $K_0$
\begin{equation}\label{sec6:eq:Dv}
\begin{aligned}
D_v&=K_0\cap U_v^{\perp}=\{z\in K_0:\ z\cdot v=z\cdot Av=0\}\\
   &=\operatorname{span}\{e_6,e_7,e_8\},
\end{aligned}
\end{equation}
$\mathcal{Q}_v$ is three-dimensional, the fibre is $\mathbb{P}^{2}$, and the admissible row spaces are
$\row\mathcal{L}=\operatorname{span}\{e_6,e_7,e_8\}\oplus\langle x\rangle$ with
$0\neq x\in\mathcal{Q}_v$. Writing $\mathcal{L}=[\ell_1;\ell_2;\ell_3;x]$ with $\ell_1,\ell_2,\ell_3$ a
basis of $D_v$, we form $H=LN$ and $H_A=LAN$; the maximal minors $m_I(x,\lambda)$ are the
$\binom{4}{3}\binom{4}{3}=16$ minors of order $3$, some identically zero, each a polynomial in
$(x,\lambda)$.


\subsection{Decisions for six directions}\label{sec6:sub:six}


We test six directions of $\mathbb{P}(E_1)$; the third is the eigen direction, the other five are
non-eigen.

\begin{table*}[t]
\centering
\caption{Decisions for six projective directions of the repeated eigenspace $E_1$. The fourth column
counts the $3\times3$ maximal minors that do not vanish identically over the function field; the fifth
column is the $\lambda$-degree of the greatest common divisor of all maximal minors over that field.
Degree $0$ means that no common root exists, hence that the fibre is generically good and, by the
full-rank leading-coefficient lemma of Section~\ref{sec:decision}, good at every point of the chart.}
\label{sec6:tab:dirs}
\begin{tabular}{@{}cccccl@{}}
\toprule
Direction $(c_1\!:\!c_2)$ & Type & Fibre & Nonzero maximal minors & $\deg_\lambda\gcd$ & Verdict\\
\midrule
$(1\!:\!0)$ & non-eigen & $\mathbb{P}^{2}$ & $2$ & $2$ & entire fibre infeasible\\
$(0\!:\!1)$ & non-eigen & $\mathbb{P}^{2}$ & $2$ & $2$ & entire fibre infeasible\\
$(1\!:\!1)$ & \textbf{eigen} & $\mathbb{P}^{3}$ & $8$ & $\mathbf{0}$ & \textbf{a feasible point exists}\\
$(1\!:\!2)$ & non-eigen & $\mathbb{P}^{2}$ & $4$ & $2$ & entire fibre infeasible\\
$(2\!:\!-1)$ & non-eigen & $\mathbb{P}^{2}$ & $4$ & $2$ & entire fibre infeasible\\
$(3\!:\!1)$ & non-eigen & $\mathbb{P}^{2}$ & $4$ & $2$ & entire fibre infeasible\\
\bottomrule
\end{tabular}
\end{table*}

\textbf{Observation.} Inside one and the same two-dimensional eigenspace, the fibre over $(1\!:\!1)$
contains a feasible row space, whereas every fibre over the other five directions is infeasible, so
feasibility on $\mathbb{P}(E_1)$ is \emph{not constant}. The differing numbers of nonzero minors ($2$,
$4$ and $8$) are a consequence of the particular directions and do not affect the decision, which uses
the gcd of \emph{all} minors.

\textbf{Why the eigen direction is the good one.} The two phenomena are related but not equivalent: the
eigen direction is the only one whose fibre is a $\mathbb{P}^{3}$ instead of a $\mathbb{P}^{2}$, so its
parameter space is one dimension larger and has more room to avoid the bad locus. This instance
therefore also illustrates the general statement that the eigen and non-eigen strata must be handled
separately; it does not assert that an eigen direction is always the feasible one.


\subsubsection{Location of the common roots}\label{sec6:sub:roots}


The common factors of the infeasible directions have the $\lambda$-parts of Table~\ref{sec6:tab:roots}:
for $(1\!:\!0)$ the polynomial $(\lambda-1)(\lambda-3)$, and a power of $(\lambda-1)$ for the other
directions, so all common roots lie in the spectrum $\{1,2,3\}$ of the compression.

\begin{table}[t]
\centering
\caption{Location of the common roots. All of them lie in the spectrum $\{1,2,3\}$ of the compression.}
\label{sec6:tab:roots}
\begin{tabular}{@{}lll@{}}
\toprule
Direction & $\lambda$-part of the common factor & Common roots\\
\midrule
$(1\!:\!0)$ & $(\lambda-1)(\lambda-3)$ & $\lambda=1,\ 3$\\
direction of $\mu=2$ & $(\lambda-1)$ & $\lambda=1$\\
direction of $\mu=3$ & $(\lambda-1)^{2}$ & $\lambda=1$\\
\bottomrule
\end{tabular}
\end{table}

\textbf{Range of the computation.} The direction space $\mathbb{P}(E_1)$ is one-dimensional and we
have taken six of its points, including the unique eigen direction; this establishes non-constancy but
is \emph{not} a classification of all directions of $\mathbb{P}(E_1)$.


\subsection{Certificates}\label{sec6:sub:cert}


\textbf{Infeasible directions.} For each of the five non-eigen directions of Table~\ref{sec6:tab:dirs},
over the function field $\mathbb{Q}(x_1,\dots,x_{r(v)-1})$ the exact computation gives
$\deg_\lambda g=2$ and $g\mid m_I$ for every maximal minor $I$, divisibility being verified by division
with remainder over that field (the remainder is zero). For $(1\!:\!0)$, with homogeneous coordinates
$(x_0\!:\!x_1\!:\!x_2)$ on the fibre, the gcd over $\mathbb{Q}(x)[\lambda]$ is defined only up to a unit
of that ring, and $x_2$ is such a unit; the representative used below is
\begin{equation}\label{sec6:eq:gexplicit}
g=x_2\,(\lambda-1)(\lambda-3),
\end{equation}
so the statement is that the gcd equals this $g$ \emph{up to a unit}; its normalized representative, monic
in $\lambda$, is the $\lambda$-part $(\lambda-1)(\lambda-3)=g/x_2$. The factor $x_2$ comes from the
parameters, the locus $x_2=0$ being the degenerate locus where $\rank H<\omega$, which is not part of the
chart.

\textbf{What the fibre computation covers: the affine cone, not one chart.} The fibre over a direction
is a projective space $\mathbb P^{\tilde r}$, and the computation is \emph{not} performed on a single
affine chart of it. The parameter is a nonzero vector $x\in\mathbb R^{\tilde r}$, $\tilde r=r(v)$, that is,
the affine cone with the origin removed; projective points are the equivalence classes
$[x]=\{tx:\ t\ne0\}$. Each maximal minor $m_I(x,\lambda)$ is homogeneous of degree at most \$ in \$re is $\lambda$ with $m_I(x,\lambda)=0$ for all $I$'' is invariant under
$x\mapsto tx$ and defines a property of $[x]$ alone. The gcd is taken in $K(x_1,\dots,x_{\tilde r})[\lambda]$,
the rational function field of the cone; $\deg_\lambda\gcd=0$ there says that no common root exists at
a \emph{general} cone point, and the full-rank lemma upgrades this to \emph{every} cone point. Hence the
conclusion concerns the whole projective fibre, as claimed, and not one chart or one sample point.

\textbf{Three assertions, and which certificate supports which.} One must distinguish
(a) \emph{some point of the fibre is feasible}; (b) \emph{a dense open subset of it is feasible};
(c) \emph{the whole fibre is infeasible}. A single point with $\gcd=1$ establishes (a); openness of the
good set then gives (b); a positive-degree gcd together with the full-rank leading-coefficient lemma
establishes (c). The counterexample uses (a) for the direction $(1\!:\!1)$ and (c) for the five
non-eigen directions; the contrast between (a) and (c) \emph{is} the non-constancy of feasibility on
$E_1$.

\textbf{Why divisibility suffices for the whole fibre.} For a general family of polynomials a positive
degree of the gcd over the function field does \emph{not} imply that every specialization has a common
root, since the common factor may specialize to a nonzero constant. What licenses the conclusion is the
full-rank leading-coefficient lemma of Section~\ref{sec:decision}: at every point of the cone
$\rank H=\omega$, hence some maximal minor has $\lambda$-degree exactly $\omega$, so a nonconstant $g$
cannot specialize to a nonzero constant; every specialization therefore carries a nontrivial common
factor. Concretely, $g\mid m_I$ is the family of divisions
\begin{equation}\label{sec6:eq:division}
m_I=g\cdot q_I+r_I,\qquad \deg_\lambda r_I<\deg_\lambda g,
\end{equation}
all with $r_I=0$, verified for each maximal minor separately in exact rational arithmetic; the
quotients are of moderate size and the expanded identities are collected in Appendix~\ref{app:data}.
The computation is carried out in the polynomial ring $\mathbb{Q}(x_0,x_1,x_2)[\lambda]$ with
$\lambda$-degree at most $3$, so neither the gcd nor the divisions require any algebraic extension.

\textbf{The feasible direction.} For the eigen direction $[1\!:\!1]$, at the rational point
\begin{equation}\label{sec6:eq:x0}
\xi_0=(1,2,3,4)
\end{equation}
of the cone over the fibre $\mathbb{P}^{3}$ the greatest common divisor of all maximal minors is
$\gcd_I m_I(\xi_0,\lambda)=1$; Appendix~\ref{app:data} gives the specialized minors and an explicit
Bezout combination $\sum_I u_I(\lambda)\,m_I(\xi_0,\lambda)=1$. The chart construction of
Section~\ref{sec:decision} turns $\xi_0$ into a concrete row space, which is then a feasible design: it
satisfies both structural conditions and the spectral condition. A single point suffices for a good
direction by (a)$\Rightarrow$(b) above; for a bad direction the whole fibre must be covered, which is
(c).


\subsection{What the counterexample does and does not show}\label{app:example:wording}



The purpose of this step is to state exactly what the instance refutes; it is used in
Corollary~\ref{sec6:cor:reduction}.

\begin{claim}[counterexample]\label{sec6:prop:nonconstant}
There is a rational instance with $n=8$, $p=2$ and $q=4$ in which the compression has a
two-dimensional real eigenspace $E_1$; inside $E_1$ the fibre over the direction $[1\!:\!1]$ contains
a feasible observer, whereas the fibres over the five non-eigen directions tested are all infeasible.
Hence feasibility is not constant on that eigenspace.
\end{claim}

\begin{cor}[failure of the representative-direction reduction]\label{sec6:cor:reduction}
The procedure ``take one eigenvector for each distinct real eigenvalue of the compression and decide
feasibility from that representative'' is not correct: the representative $(1\!:\!0)$ would declare the
instance of Claim~\ref{sec6:prop:nonconstant} infeasible, whereas the instance is feasible.
\end{cor}

The counterexample shows that the direction parameter must retain the \emph{whole} projective eigenspace
$\mathbb{P}(E_\mu)$: the complexity theorem of Section~\ref{sec:complexity} takes geometric simplicity of
the real spectrum as a hypothesis, and the general decision problem for repeated eigenspaces lies
outside its scope; the direction is also a genuine degree of freedom which cannot be absorbed into the
row-space parameters.


The limitations of this conclusion are collected in Section~\ref{sec6:sub:conclusion}; what is shown
here is only the failure of one reduction, and the untested directions of $E_1$ are left alone.

\subsection{Complete data and exact certificates}\label{appA:data}

The following gives the complete input and the exact certificates for the counterexample above; all
arithmetic is over the rationals, with no rounding.

\subsubsection{Input and derived quantities}


Take $n=8$, $p=2$, $q=4$ rows and $L_0$ empty ($r_0=0$), with
\begin{equation}
\label{appA:eq:C}
C=\begin{bmatrix}1&0&0&0&0&0&0&0\\0&1&0&0&0&0&0&0\end{bmatrix},
\end{equation}
\begin{equation}
\label{appA:eq:A}
A=\begin{bmatrix}
0&0&1&0&0&0&0&0\\0&0&0&1&0&0&0&0\\
1&-2&1&0&-1&1&1&-1\\-1&2&1&0&2&-2&-1&2\\
0&0&0&0&1&0&0&0\\0&0&0&0&0&1&0&0\\
0&0&0&0&0&0&2&0\\0&0&0&0&0&0&0&3
\end{bmatrix}.
\end{equation}
For the data of \eqref{appA:eq:C},
\begin{equation}\label{appA:eq:rank}
\rank\begin{bmatrix}C\\ CA\end{bmatrix}=4,\qquad k=4,\qquad \rho=4,\qquad \omega=3,
\end{equation}
and a basis of $K_0$, the compression, the characteristic polynomial and the complementary block are
\begin{gather}
\label{appA:eq:N}
N=\begin{bmatrix}
0&0&0&0\\0&0&0&0\\0&0&0&0\\0&0&0&0\\
1&0&0&0\\0&1&0&0\\0&0&1&0\\0&0&0&1
\end{bmatrix},\\
M=(N^{\mathsf T}N)^{-1}N^{\mathsf T}AN=\operatorname{diag}(1,1,2,3),\\
\label{appA:eq:chi}
\chi_M(t)=(t-3)(t-2)(t-1)^2,\\
A_{21}=\begin{bmatrix}
0&0&0&0\\0&0&0&0\\-1&1&1&-1\\2&-2&-1&2\\
0&0&0&0\\0&0&0&0\\0&0&0&0\\0&0&0&0
\end{bmatrix}.
\end{gather}
Two rows of $A_{21}$ are nonzero; this is the map that distinguishes eigen from non-eigen directions.


\subsubsection{The six directions}


The eigenspace of $\mu=1$ is $E_1=\operatorname{span}\{e_1,e_2\}\subseteq\mathbb R^4$ in $K_0$
coordinates, so $\mathbb P(E_1)\cong\mathbb P^1$; $(1,1)$ is the unique eigen direction
($A_{21}(c_1,c_2,0,0)^{\mathsf T}=0\iff c_1=c_2$), and its fibre $\mathbb P^3$ has one parameter more
than the fibres $\mathbb P^2$ of the other directions. The six directions and their verdicts are listed
in Table~\ref{appA:tab:dirs}.


\subsubsection{Chart variables, and the certificates in expanded form}


For a non-eigen direction the fibre is $\mathbb P^2$ with homogeneous coordinates $(x_0\!:\!x_1\!:\!x_2)$;
for the eigen direction $(1\!:\!1)$ it is $\mathbb P^3$ with $(x_0\!:\!x_1\!:\!x_2\!:\!x_3)$. Below,
every maximal minor is written as a polynomial in $\lambda$ whose coefficients are polynomials in those
coordinates with integer coefficients, so that each identity can be checked by expansion. All
coefficients lie in $\mathbb Q$.

\emph{Bad directions.} In each of the five non-eigen directions the maximal minors share the factor
\begin{equation}\label{appA:eq:g}
g(x,\lambda)=x_2(\lambda-1)(\lambda-3)=x_2\lambda^2-4x_2\lambda+3x_2 .
\end{equation}
With $\deg_\lambda g=2=\deg_\lambda\gcd$ and $g\mid m_I$ for every $I$, the gcd is $g$ up to a unit of
$\mathbb Q(x)[\lambda]$, and $x_2$ is such a unit; the normalized (monic in $\lambda$) representative is
$(\lambda-1)(\lambda-3)=g/x_2$, and the $\lambda$-degree, which is what the decision uses, is the same for
both normalizations.
For $(1\!:\!0)$ the two nonzero maximal minors are, explicitly,
\begin{align}
m_1(x,\lambda)&=(\lambda-1)(\lambda-2)(\lambda-3)=\lambda^3-6\lambda^2+11\lambda-6,
\label{appA:eq:m1}\\
m_2(x,\lambda)&=x_2(\lambda-1)(\lambda-3)=x_2\lambda^2-4x_2\lambda+3x_2,
\label{appA:eq:m2}
\end{align}
and division in $\mathbb Q(x)[\lambda]$ gives, with zero remainders in both cases,
\begin{equation}\label{appA:eq:divcert}
m_1=g\cdot\frac{\lambda-2}{x_2},\qquad m_2=g\cdot 1 .
\end{equation}
For the remaining non-eigen directions the same factor $g$ occurs; the counts and the outcome of the
divisions are collected in Table~\ref{appA:tab:cert}.

\begin{table}[t]
\centering
\caption{Divisibility certificate per direction. ``Minors'' is the number of maximal minors that are not
identically zero; ``division'' records whether $g\mid m_I$ with zero remainder for every one of them.}
\label{appA:tab:cert}\label{appA:tab:dirs}
\begin{tabular}{@{}c c c c l@{}}
\toprule
Direction & Minors & $\deg_\lambda g$ & Division & Verdict \\
\midrule
$(1,0)$  & 2 & 2 & yes & entirely bad \\
$(0,1)$  & 2 & 2 & yes & entirely bad \\
$(1,1)$  & 8 & 0 & --- ($\gcd=1$) & good at $\xi_0$ \\
$(1,2)$  & 4 & 2 & yes & entirely bad \\
$(2,-1)$ & 4 & 2 & yes & entirely bad \\
$(3,1)$  & 4 & 2 & yes & entirely bad \\
\bottomrule
\end{tabular}
\end{table}

\emph{The good direction.} At the rational point
\begin{equation}\label{appA:eq:xipoint}
\xi_0=(1,2,3,4)\in\mathbb P^3
\end{equation}
of the fibre over $(1\!:\!1)$ the eight nonzero maximal minors specialize to
\begin{gather}
-(\lambda-1)(\lambda-2)(\lambda-3),\quad +(\lambda-1)(\lambda-2)(\lambda-3),\notag\\
5(\lambda-1)(\lambda-2),\quad -5(\lambda-1)(\lambda-2),\notag\\
(\lambda-1)(\lambda-3),\quad -(\lambda-1)(\lambda-3),\notag\\
-5(\lambda-2)(\lambda-3),\quad 5(\lambda-2)(\lambda-3),\label{appA:eq:spec}
\end{gather}
in the order $m_1,\dots,m_8$, and the Bezout identity
\begin{equation}\label{appA:eq:bezout}
\frac{\lambda-4}{10}\,m_4+\frac{\lambda-4}{2}\,m_5-\frac{1}{10}\,m_7=1
\end{equation}
holds. Its left-hand side is
$\frac{1}{2}\bigl[-(\lambda-4)(\lambda-1)(\lambda-2)+(\lambda-4)(\lambda-1)(\lambda-3)+(\lambda-2)(\lambda-3)\bigr]
=\frac{1}{2}\bigl[-(\lambda^2-5\lambda+4)+(\lambda^2-5\lambda+6)\bigr]=1$,
so \eqref{appA:eq:bezout} is an identity in $\mathbb Q[\lambda]$. Consequently
$\gcd_I m_I(\xi_0,\lambda)=1$, and $\xi_0$ is a feasible point of that fibre.

\emph{What these certificates establish.} Identities \eqref{appA:eq:divcert} and
\eqref{appA:eq:bezout} are finite rational identities: expanding either side and cancelling verifies
them in any exact arithmetic system, with no algebraic extension and no rounding. Together with the
full-rank leading-coefficient lemma of Section~\ref{sec:decision} they give the two conclusions used in
Appendix~\ref{app:example}: the fibres over the five non-eigen directions are entirely bad, and
the fibre over $(1\!:\!1)$ contains a feasible design.


\subsubsection{Reproduction}


Compute $N$, $M$, $\chi_M$, $A_{21}$ from $A$ and $C$ by rational elimination; for each direction build
$D_v$, $\mathcal Q_v$ and the maximal minors of $P(\lambda)=\lambda H-H_A$; divide them by $g$ in
$\mathbb Q(x)[\lambda]$ (fraction-free subresultants) to reproduce Table~\ref{appA:tab:cert}; substitute
$\xi_0$ and expand to verify \eqref{appA:eq:spec} and \eqref{appA:eq:bezout}.


\normalsize



\begin{thebibliography}{99}
\bibitem{fer26a} J. A. Fernando. A minimum-order functional observer beyond the Darouach and Luenberger constructions. arXiv:2609.33615, 2026 (preprint).
\bibitem{fer26b} J. A. Fernando. Complete characterization of minimum-order functional observers from Darouach to Luenberger. arXiv:2609.37992, 2026 (preprint).
\bibitem{fer10} J. A. Fernando, H. Trinh, L. Jennings. Functional observability and the design of minimum order linear functional observers. \emph{IEEE Trans. Automat. Control} 55(5):1268--1273, 2010.
\bibitem{dar25} M. Darouach, J. A. Fernando. On functional observability and functional observer design. \emph{Automatica}, vol. 173, Art. no. 112115, 2025.
\bibitem{hj13} R. A. Horn, C. R. Johnson. \emph{Matrix Analysis}. 2nd ed. Cambridge University Press, 2013.
\bibitem{bpr06} S. Basu, R. Pollack, M.-F. Roy. \emph{Algorithms in Real Algebraic Geometry}, 2nd ed. Berlin, Germany: Springer, 2006, Algorithms and Computation in Mathematics, vol. 10.
\bibitem{lan02} S. Lang. \emph{Algebra}. Revised 3rd ed. Graduate Texts in Mathematics 211, Springer, 2002.
\bibitem{vzgg03} J. von zur Gathen, J. Gerhard. \emph{Modern Computer Algebra}. 2nd ed. Cambridge University Press, 2003.
\bibitem{companion1} Deciding functional observer feasibility: compression, certificates, and parameterized algorithms. Companion manuscript, 2026.
\bibitem{companion2} Minimum order and minimum-order construction for functional observers. Companion manuscript, 2026.
\end{thebibliography}

\end{document}
